\documentclass[12pt, a4paper]{article}
\usepackage[utf8]{inputenc}
\usepackage{amsmath, amssymb, amsthm, amsfonts,mathrsfs}
\usepackage{CJKutf8}
\usepackage[margin=1in]{geometry}
\usepackage{tikz-cd}
\newcommand{\spec}{\operatorname{Spec}}
\newcommand{\spa}{\operatorname{Span}}
\newcommand{\p}{\mathfrak{p}}
\renewcommand{\a}{\mathfrak{a}}
\renewcommand{\b}{\mathfrak{b}}
\newcommand{\z}{\mathbb{Z}}
\newcommand{\Q}{\mathbb{Q}}
\newcommand{\R}{\mathbb{R}}
\newcommand{\q}{\mathfrak{q}}
\newcommand{\tr}{\text{Tr}}
\renewcommand{\o}{\mathcal{O}}
\newtheorem{lemma}{Lemma}


\newenvironment{solution}
  {\paragraph{\textit{Solution:}}\par\medskip}
  {\hfill$\blacksquare$\vspace{0.1cm}}

\title{Algebraic Number Theory:Assignment}
\author{
    \begin{CJK*}{UTF8}{gbsn}
    钟皓宇 (12533556)
    \end{CJK*}
}
\date{\today}

\begin{document}

\maketitle
\newcounter{paranum} 
\newcommand{\np}{\stepcounter{paranum}\noindent\textbf{\arabic{paranum}.}\quad}

\section*{Some Examples}
Let $R$ be an integral domian with fraction field $F$. Let $K$ be a field extension $F$ and let $V$ be a finite dimensional $K$-vector space.

\np $L$ is an $R$-submodule of $F\otimes_{R} L$. To anlysis, we can consider 
\begin{align*}
    (v_1,\dots,v_m) = (u_1,\dots,u_n)P
\end{align*} 
Where $(u_i)$ and $(v_i)$ is the basis of $F\otimes_{R} L$ and generators of $L$.
We can also extend $v_1,\dots,v_m$ to a basis of $F\otimes_{R} L$ by some consideration.

\np An example of $R$-lattice not support in $V$. 
Let $K=V=\mathbb{R}$ and $R=\z$. Then
$L= \z+\z\sqrt{2}$ is a $\z$-lattice but not support in $\mathbb{R}$.

\np It is easy to see that from 2., $\dim L\ge \dim_K\spa_K{L}$ and a system of vectors in $L$ is linearly independent over $R$ iff. over $K$.

\np We can clearly see from 2. the equivalence of existence of $K$-free basis and support in $V$.

\np Let $L= 2\cdot \z+\z(1+\sqrt{-5})$, $K = \Q(\sqrt{2})$ and $V = \Q(\sqrt{2},\sqrt{-5})$. Then $L$ is support in $V$.
$L$ is a non-trivial generator of ideal class group of $\Q(\sqrt{-5})$.


\np Let $L= \z+\z\sqrt{-5}$ and $L' = 2\cdot\z+\z(1+\sqrt{-5})$. Then $\det(L/L')=2 \cdot \z^*$ and $2L\subseteq L'$.

\textbf{Norm and discriminant in Algebraic Number Theory}

We consider $L/K$ is a seperable finite algebraic extension. We specifically let $O_K$ is not a PID. 
Let $I$ be a fraction ideal of $O_L$. Then there exists a basis $\{\omega_i\}$ of $L$ over $K$ such that
$$ I = \mathfrak{a}_1 \omega_1 \oplus \mathfrak{a}_2 \omega_2 \oplus \dots \oplus \mathfrak{a}_n \omega_n. $$
Similarly, we have 
$$ \mathcal{O}_L = \mathfrak{b}_1 \nu_1 \oplus \mathfrak{b}_2 \nu_2 \oplus \dots \oplus \mathfrak{b}_n \nu_n.$$
Suppose that $P$ is the transform matrix such that $(\omega_i) = P(\nu_i)$.
In this sense, the norm of $I$ is given by $$\det(P)\cdot \frac{\mathfrak{a}_1\cdots \mathfrak{a}_n}{\mathfrak{b}_1\cdots \mathfrak{b}_n}$$
The discriminant of $I$ is given by 
$$(\mathfrak{a}_1 \mathfrak{a}_2 \dots \mathfrak{a}_n)^2 \cdot \det\left(\text{Tr}_{L/K}(\omega_i \omega_j)\right).$$
The discriminant of field $L$ is the discriminant of $O_L$.
If $L=K(\theta)$ with miminal polynomial $p(x)$ then the discriminant $\text{dics}(f)$ of $f$ is given by the discriminant of lattice generated by $\{\theta^n\}$.

\np Let $L= \z+\z\sqrt{-5}$ and $L' = \z(3-\sqrt{-5})+\z(9+\sqrt{-5})$. Let $P$ be the transition matrix. Here is the triangular process of $P$.
Let $L'^{(1)}= 3\mathbb{Z}$ and $L'^{(2)} = L$. Then $\pi_1(L'^{(1)}) = 3\z$ and $\pi_2(L'^{(2)}) = \sqrt{-5}\z$ is principal.
Let $d=\det(L/L') = 12$. Then $12\mathbb{Z}\subseteq 3\mathbb{Z}$ and $12 \z\sqrt{-5}\subseteq \z\sqrt{-5}$ it implies that $\pi_i(L'^{(i)})$ is not zero.
Choose $f_i\in L'^{(i)}$ such that $\pi_i(f_i)=t_{ii}$ where $t_{ii}$ is the generator of ideal $\pi_i(L'^{(i)})$.

The core of process above is the coefficient ideal $\pi_i(L'^{(i)})$. And the condition $L'\subseteq L$ guarantee the projection is not zero ideal.

\np Resay a trick. If we don't have $L'\subseteq L$. we just need to consider $dL'\subseteq L$.

We consider field extension $L/K$ equipped with bilinear form $\tr_{L/K}$.

\np If $L/K$ is separable then $\tr$ is not singular.

\np If $L/K$ is inseparable extension, then $\tr$ is totally singular.
Let $L=\mathbb{F}_p(t^{1/p})$ and $K = \mathbb{F}_p(t)$ then $\tr(L)=p$.

\np Isotropic case: Let $L=\Q(i)$ and $L=\Q$. Consider $\tr((1+i)^2)= 0$.

\np Anisotropic case: Let $L=\Q(\sqrt{2})$ and $L=\Q$. Consider $\tr(x^2)= x^2+\bar{x}^2\neq 0$ if $x\neq 0$.

\np Alternative case: Twisted Trace. Let $L=\Q(i)$ and $L=\Q$ with $B(x,y)=\tr(i\cdot \bar{x}y)$.


\textbf{Twisted Trace Form and Different}

Let $L/K$ be a finite seperable algebraic extension. The Different of $L/K$ is defined by $(^* \o_L)^{-1}$.

\np Conisder vector space $\mathbb{F}_2x$ with quadratic form $Q(x)= 1$. Then radical is $\mathbb{F}_2x$ but it is not singular.




\section*{A.1 \quad Homework 1}

Throughout what follows, let $R$ denote an integral domain with fraction field $K$.

\subsection*{Exercise A.1.1.} 
Suppose $R$ is a local domain. Let $I$ be a fractional ideal of $R$. Prove that $I$ is invertible if and only if it is principal. \\
\emph{Hint: There exist $x \in I$, $y \in I^{-1}$ such that $xy \in R^{\times}$. Then prove $I = xR$.}
\begin{solution}
    Suppose $x\neq 0$ and $I=xR$ is principal then let $J=x^{-1}R$ and $IJ=R$ implies $I$ is invertible. Conversely, if there exist a fractional ideal $J$ such that $IJ=R$. Then there exist $x_i\in I$ and $y_i\in J$ such that
    \begin{align*}
        \sum x_iy_i = 1.
    \end{align*} 
    Because $R$ is local, without loss of generality, we can suppose $x_1y_1\in R\backslash \mathfrak{m}=R^{\times}$. Since $x_1\in I$, $x_1R\subseteq I$ is obvious. Hence $y_1\in J$ implies $y_1I\subseteq R$, multiplying by $x_1$ on both sides, we obtain that 
    $I\subseteq x_1 R$. Therefore, $I=x_1R$ is principal.
\end{solution}


\subsection*{Exercise A.1.2.} 
Prove that a Dedekind domain with only finitely many prime ideals is a PID.
\begin{solution}
    Let $\mathfrak{p}_1,\dots,\mathfrak{p}_n$ be all prime ideals of Dedekind domain $R$.
    It suffices to prove all prime ideals are principal. Without loss of generality, we prove $\mathfrak{p}_1$ is principal and others are in 
    the same way. Let $x\in \mathfrak{p}_1\backslash \mathfrak{p}_1^2$ and by Chinese Remainder Theorem, there exist $u\in R$ such that 
    \begin{align*}
        u\equiv x \bmod \mathfrak{p}_1^2 \text{ and } u\equiv 1 \bmod \mathfrak{p}_i\text{ for }i\neq 1.
    \end{align*} 
    Let $v_1,\dots,v_n$ be the valuations of local rings $R_{\mathfrak{p}_1},\dots,R_{\mathfrak{p}_n}$ respectively. Then we obtain that 
    \begin{align*}
        v_1(u) = 1 \text{ and } v_i(u) = 0\text{ for }i\neq 1,
    \end{align*}
    because $u-x\in \mathfrak{p}_1^2$ and implies $u\in (x)+\mathfrak{p}_1^2\subseteq \mathfrak{p}_1$ meaning that $v_1(u)\ge 1$. Suppose $v_1(u)\ge 2$, equivalently $u\in \mathfrak{p}_1^2$, then $x\in (u)+\mathfrak{p}_1^2 = \mathfrak{p}_1^2$ contradicting our
    assumption, therefore $v_1(u)=1$. Because $R$ is Dedekind domain, the unique factorization of fractional ideals implies 
    \begin{align*}
        (u) = \mathfrak{p}_1^{v_1(u)}\dots\mathfrak{p}_n^{v_n(u)} = \mathfrak{p}_1,
    \end{align*}
    the principality of $\mathfrak{p}_1$ follows.
\end{solution}



\subsection*{Exercise A.1.3.} 
Suppose $R$ is a DVR with maximal ideal $\mathfrak{p}$. Let $k = R/\mathfrak{p}$ and $U = R^{\times}$. Let $v$ be the normalized discrete valuation of $K$ associated to $R$. For each $n \geq 1$, define
\[
U_n = 1 + \mathfrak{p}^n = \{x \in K \mid v(x - 1) \geq n\} \, .
\]

\begin{enumerate}
    \item Show that each $U_n$ is a subgroup of $U$ and that $\bigcap_{n=1}^{\infty} U_n = 1$.
    \item Prove that the natural quotient map $R \to k = R/\mathfrak{p}$ induces an isomorphism of multiplicative groups $U/U_1 \cong k^*$.
    \item Prove that for each $n \in \mathbb{Z}$, there is an isomorphism of $R$-modules $k = R/\mathfrak{p} \xrightarrow{\sim} \mathfrak{p}^n/\mathfrak{p}^{n+1}$.
    \item Prove that for each $n \geq 1$, the map $u \mapsto u - 1$ induces an isomorphism of abelian groups $U_n/U_{n+1} \xrightarrow{\sim} \mathfrak{p}^n/\mathfrak{p}^{n+1}$.
    \item Let $p$ be the characteristic of $k$. (So $p = 0$ or $p$ is a prime number.)
    \begin{enumerate}
        \item[(a)] Show that $U_n^p := \{a^p \mid a \in U_n\}$ is contained in $U_{n+1}$ for every $n \geq 1$.
        \item[(b)] Let $m \geq 1$ be an integer not divisible by $p$. Prove that for every $n \geq 1$, the group homomorphism
        \[
        U_n \longrightarrow U_n \ ; \quad u \longmapsto u^m
        \]
        is injective.
    \end{enumerate}
\end{enumerate}

\begin{solution}
    \textbf{(1)} Closed under multiplication holds since $(1 + \mathfrak{p}^n)(1 + \mathfrak{p}^n)\subseteq 1 + \mathfrak{p}^n$. 
    Identity element is $1$. Suppose $x\in U_n$, then $v(x^{-1}-1) = v(x^{-1})+v(1-x)$ hence $x$ is invertible meaning that $v(x)=0$ and 
    $v(1-x) = v(x-1)$ implies $v(x^{-1}-1)\ge n$. Suppose $x\in \bigcap_{n=1}^{\infty} U_n $, then we obtain $v(x-1) = \infty$, equivalently, $x=1$.
    \textbf{(2)} The natural quotient map restricted to $U$, denoted by $\pi:U\to k^*$, induces isomorphism $U/U_1 \cong k^*$ because $\ker(\pi) = U_1$.
    \textbf{(3)} Consider commutative graph 
\[
\begin{tikzcd}[row sep=large, column sep=large]
0 \arrow[r] & \mathfrak{p} \arrow[r, hook] \arrow[d, "\cdot u^n"', "\cong"] & R \arrow[r, two heads] \arrow[d, "\cdot u^n"', "\cong"] & k \arrow[r] \arrow[d, dashed, "\bar{\phi}", "\cong"'] & 0 \\
0 \arrow[r] & \mathfrak{p}^{n+1} \arrow[r, hook] & \mathfrak{p}^n \arrow[r, two heads] & \mathfrak{p}^n/\mathfrak{p}^{n+1} \arrow[r] & 0
\end{tikzcd}
\] 
which the third isomorphism follows from the Five Lemma.
\textbf{(4)} Consider a surjective group homomorphism 
\begin{align*}
    \varphi:U_n\to \mathfrak{p}^n/\mathfrak{p}^{n+1} \quad u\mapsto u-1 + \mathfrak{p}^{n+1}
\end{align*}
with kernel $U_{n+1}$, then the first isomorphism theorem yields isomorphism $U_n/U_{n+1} \cong \mathfrak{p}^n/\mathfrak{p}^{n+1}$.
\textbf{(5.a)} For $p=0$, obvious to see that $U_n^0=\{1\}\subseteq U_{n+1}$. For $p\neq 0$ and any $a\in U_n$, 
\begin{align*}
    v(a^p-1) =  v((a-1)^p) = p\cdot v(a-1)\ge n+1
\end{align*} 
implies $a^p\in U_{n+1}$.
\textbf{(5.b)} It suffices to prove for $m=q$ where $q$ is a prime number and coprime to $p$. Suppose that $u\in U_n$ such that $u^q =1$. 
According to Bézout's identity and \textbf{(5.a)}, suppose $ap+bq=1$, then $u = u^{ap+bq} = u^{ap}\in U_{n+1}$. Thus we can consider a new group homomorphism:
$U_{n+1} \to U_{n+1}$. Continuing this process, we obtain that $u\in \bigcap_{n=1}^{\infty} U_n$, then $u=1$ and injectivity holds.
\end{solution}

\subsection*{Exercise A.1.4.} 
Prove that if a domain $R$ is integrally closed, then so is the polynomial ring $R[t]$.
\begin{solution}
Let $f=p/q$ be an arbitrary integral element in $R(t)$. Without loss of generality, we can suppose $q$ is irreducible, because if $q=q_1q_2$ then we consider
integral element $fq_1$, and suppose $p,q$ are coprime. Suppose $f$ satisfies a monic polynomial $f^n+a_{n-1}f^{n-1}+\dots+a_0=0$ where $a_k\in R[t]$. Then, by multiplying both side by $q^n$ 
, we obtain that $p^n+a_{n-1}qp^{n-1}+\dots+a_0q^n=0$, which implies that $q|p^n$. Let $F=\text{Frac}(R)$ be the fraction field of $R$, and 
interprete $p,q$ as elements in $F[t]$, then we obtain that $q|p$ since $F[t]$ is an UFD and hence $q\in F^*$ because $p,q$ coprime.
Therefore, we can write $f=c_mt^m+\dots+c_0$ where $c_k\in F$.  Let $f_1 = t^r - f$, where $r=\deg f+\sum \deg a_k$, then $f_1$ is the root of polynomial
\begin{align*}
    P(X)=(t^r-X)^{n}+a_{n-1}(t^r-X)^{n-1}+\dots + a_0.
\end{align*}
Let $X=0$, we obtain the constant term of this polynomial
\begin{align*}
    t^{rn}+a_{n-1}t^{rn-r}+a_{n-2}t^{rn-2r}+\dots+a_0
\end{align*}
denoted by $P(0)$. Let $P(X)=XQ(X)+P(0)$, then from $P(f_1)=0$ we obtain that 
\begin{align*}
    P(0) = -f_1Q(f_1),
\end{align*} 
the roots of $f_1$ are among the roots of monic polynomial $P(0)\in R[t]$, the coefficients of $f$ as of $f_1$ are elementary symmetric polynomials of its roots, which are integral over $R$. Therefore, we obtain that the coefficients of 
$f$ are contained in $R$, $R[t]$ is integrally closed.  
\end{solution}



\subsection*{Exercise A.1.5.} 
Let $A = \mathbb{Q}[X, Y]/(X^2 - Y^3)$. Show that $A$ is a domain but it is not integrally closed.
\begin{solution}
    $(X^2-Y^3)$ is a prime ideal because $X^2-Y^3$ is irreducible and $\mathbb{Q}[X,Y]$ is UFD, \textit{a fortiori}, $A$ is a domain.
    Hence, $X/Y$ is the root of a monic polynomial $P(Z) = Z^2-Y$ but not contained in $A$. If it is contained, then there exist $f,g\in\mathbb{Q}[X,Y]$ such that 
    \begin{align*}
        X = Yf-(X^2-Y^3)g
    \end{align*}
    implies contradiction by comparing their linear terms.
\end{solution}

\subsection*{Exercise A.1.6.} 
Let $R$ be a Dedekind domain and let $\mathfrak{a} \subseteq R$ be a nonzero ideal. Show that the quotient ring $R/\mathfrak{a}$ has only finitely many ideals and that all of these ideals are principal.
\begin{solution}
    By isomorphism theorem of rings, the natural quotient $R\to R/\mathfrak{a}$ gives one-to-one correspondence of the ideals of $R$ containing $\mathfrak{a}$ and the ideals of $R/\mathfrak{a}$.
    Therefore, all ideals of $R/\mathfrak{a}$ can be written as the form 
    \begin{align*}
        \mathfrak{p}_1^{i_1}\dots \mathfrak{p}_n^{i_n} + \mathfrak{a}
    \end{align*}
    where $\mathfrak{p}_k$ are the prime ideals of $R$ containing $\mathfrak{a}$ and $0\le i_k\le v_k(\mathfrak{a})$ where $v_k$ is the valuation associated to $\mathfrak{p}_k$.
    The principality directly from \textit{Exercise A.1.2.}
\end{solution}
\subsection*{Exercise A.1.7.} 
Show that every integral ideal of a Dedekind domain can be generated by two elements.
\begin{solution}
    Let $\mathfrak{a}$ is an arbitrary ideal of $R$ with prime ideal decomposition
    \begin{align*}
        \mathfrak{a} = \mathfrak{p}^k\mathfrak{b}
    \end{align*}
    Where $\mathfrak{p}$ is a prime ideal coprime to the ideal $\mathfrak{b}$.
    The quotient ring $R/\mathfrak{p}^{k+1}\mathfrak{b}$ is principal from \textit{Exercise A.1.6.} Therefore
    we can assume that $\mathfrak{a} = (a)+\mathfrak{p}^{k+1}\mathfrak{b}$. Hence, consider another quotient ring $R/(a)$, we obtain that $\mathfrak{a}=(a,b)$. 
\end{solution}

\section*{A.2 \quad Homework 2}
\subsection*{Exercise A.2.1}
Let $R$ be the ring of algebraic integers in $\mathbb{C}$, i.e., the integral closure of $\mathbb{Z}$ in $\mathbb{C}$. Consider the ideal $I \subseteq R$ generated by the infinite subset $\{\sqrt[2^n]{2} \mid n \in \mathbb{N}\}$. 

Show that $I$ is not finitely generated. Deduce that $R$ is not Noetherian (and hence not a Dedekind domain).

\begin{solution}
    Suppose that $I$ is finitely generated with generators $\alpha_1,\dots,\alpha_m$.
    According to Chevalley's Extension Theorem for Valuations, we can extend 2-adic valuation on $\mathbb{Q}$ to $\overline{\mathbb{Q}}$, let $v_2$ denote this extended valuation.
    By definition, we obtain that $R\subseteq \overline{\mathbb{Q}}$ and $v_2(R)\ge 0$.
    Since $I$ is generated by $\{\sqrt[2^n]{2} \}$ and hence $v_2(\sqrt[2^n]{2})= 2^{-n}$ is strictly greater than zero, the valuation $v_2(\alpha_i)$ is strictly greater than zero.
    Let $\varepsilon = \min\{v_{2}(\alpha_i)\}$ be the lower bound of valuations of finite generators $\alpha_i$.
    For arbitrary $r \in I$, we can write $r = \sum_{i=1}^m c_i \alpha_i$ for some $c_i \in R$. Using the non-Archimedean property and $v_2(c_i) \ge 0$, we obtain:$$v_2(r) \ge \min_{i} \{v_2(c_i) + v_2(\alpha_i)\} \ge \min_{i} \{v_2(\alpha_i)\} = \varepsilon > 0$$
    However we find that $v_{2}(\sqrt[2^n]{2}) = 2^{-n}$ converging to zero, which is a contradiction.
\end{solution}
\subsection*{Exercise A.2.2}
Let $R$ denote a Dedekind domain with fraction field $K$ and suppose $R \neq K$. 

Prove that there exists a set $\Omega$ of normalized discrete valuations of $K$ such that the following two properties hold:
\begin{itemize}
    \item[(R1)] $R = \{x \in K \mid v(x) \geq 0 \text{ for all } v \in \Omega\}$;
    \item[(R2)] for every $a \in K^\times$, the set $\{v \in \Omega \mid v(a) \neq 0\}$ is finite.
\end{itemize}

\begin{solution}
    Let 
    \begin{align*}
        \Omega = \{v_{\mathfrak{p}} \mid \mathfrak{p} \in \text{Spm}(R)\}
    \end{align*}
    where $v_{\mathfrak{p}}$ is the normalized discrete valuation associated with $R_{\mathfrak{p}}$.
    For (R1), we find that $v_{\mathfrak{p}}(x)\ge 0$ if and only if $x\in R_{\mathfrak{p}}$, then we obtain that
    \begin{align*}
        \{x \in K \mid v(x) \geq 0 \text{ for all } v \in \Omega\} = \bigcap_{\mathfrak{p}\in\text{Spm}(R)}R_{\mathfrak{p}} = R.
    \end{align*}
    For $(R2)$, we first fix an arbitrary element $a\in K^*$ and let $A$ denote the set $\{v \in \Omega \mid v(a) \neq 0\}$.
    Since $R$ is Dedekind, we may suppose that 
    \begin{align*}
        aR = \mathfrak{p}_1^{a_1}\dots \mathfrak{p}_n^{a_n}.
    \end{align*}
    where $a_i\neq 0$.
    I claim that $A = \{v_{\mathfrak{p}_1}, \dots, v_{\mathfrak{p}_n}\}$. For an arbitrary maximal ideal $\mathfrak{q} \notin \{\mathfrak{p}_1, \dots, \mathfrak{p}_n\}$, it does not appear in the factorization of $aR$, yielding $v_{\mathfrak{q}}(a) = 0$ and hence $v_{\mathfrak{q}} \notin A$. This implies $A \subseteq \{v_{\mathfrak{p}_1}, \dots, v_{\mathfrak{p}_n}\}$. Conversely, by definition we have $v_{\mathfrak{p}_i}(a) = a_i \neq 0$, meaning $\{v_{\mathfrak{p}_1}, \dots, v_{\mathfrak{p}_n}\} \subseteq A$. 
    Therefore, $A$ is precisely this finite set. 
\end{solution}

\subsection*{Exercise A.2.3}
Let $R$ denote an integral domain with fraction field $K$. Let $\Omega$ be a set of normalized discrete valuations of $K$ satisfying Properties (R1) and (R2) in Exercise A.2.2. For each $v \in \Omega$, let $R_v$ be its valuation ring and let $\mathfrak{p}_v$ be the maximal ideal of $R_v$.

Suppose further that every nonzero prime ideal of $R$ is maximal.

\begin{enumerate}
    \item Let $\mathfrak{p}$ be a maximal ideal of $R$.
    \begin{enumerate}
        \item Let $z \in K^\times$ be such that $R \cap z^{-1}R \subseteq \mathfrak{p}$. Show that there exists $v \in \Omega$ such that $\mathfrak{p} = R \cap \mathfrak{p}_v$ and $v(z) < 0$.
        \item Prove that $\Omega_{\mathfrak{p}} := \{v \in \Omega \mid R_v \supseteq R_{\mathfrak{p}}\}$ is a finite set.
        \item Prove that $R_{\mathfrak{p}} = R_v$ for a unique $v \in \Omega$. Thus $\Omega_{\mathfrak{p}}$ consists of a single element.
    \end{enumerate}
    
    \item Prove that for every $v \in \Omega$, there exists a unique maximal ideal $\mathfrak{p}$ such that $R_v = R_{\mathfrak{p}}$.
    
    \item Deduce from (1) and (2) that the map
    \[ \Omega \longrightarrow \text{Spm}(R) \ ; \ v \longmapsto \mathfrak{p}_v \cap R \]
    is well defined and bijective. Let $\mathfrak{p} \mapsto v_{\mathfrak{p}}$ denote its inverse map.
    
    \item Let $I \subseteq R$ be a nonzero ideal.
    \begin{enumerate}
        \item Prove that $\{\mathfrak{p} \in \text{Spm}(R) \mid v_{\mathfrak{p}}(I) \neq 0\}$ is finite, where $v_{\mathfrak{p}}(I) := \min\{v_{\mathfrak{p}}(x) \mid x \in I\}$.
        \item Prove that there is a unique factorization $I = \prod_{\mathfrak{p} \in \text{Spm}(R)} \mathfrak{p}^{r_{\mathfrak{p}}}$, where each $r_{\mathfrak{p}} \in \mathbb{N}$ and almost all of them are zero.
        \item If $J \subseteq R$ is another nonzero ideal, show that $I \mid J$ if and only if $I \supseteq J$.
    \end{enumerate}
    
    \item Prove that $R$ is a Dedekind domain.
\end{enumerate}
\begin{solution}
(1.a) Let $\Gamma =\{v\in \Omega|v(z)<0\}$. According to (R2), it is finite. Hence, it cannot be emptyset since if so then $z^{-1}R\cap R = R \subseteq \p$ is impossible.
We have 
\begin{align*}
    R\cap z^{-1}R = \{r\in R|v(r)\ge -v(z)\text{ for any }v\in \Gamma\}.
\end{align*}
If our claim is not hold, that is, for any $v\in \Gamma$ we then have $\p\neq R\cap \p_v$.
We obtain pairwise coprime ideals $\p,\p_i$ (Here we can suppose that $\p_i\cap R\neq \p_j\cap R$ by discarding the same ideals and choose the maximal valuation on $z$)
and consider CRT by 
\begin{align*}
    r \equiv 1 \bmod \p \quad r\equiv  0 \bmod (\p_{v_i}\cap R)^{-v_i(z)}
\end{align*}
Then we find that $r\in R\cap z^{-1}R$ and $r\notin \p$. It is a contradiction.

(1.b)
Suppose that $\Omega_{\p}$ is an infinite set and for arbitrary $v\in\Omega_\p$ we can have
\begin{align*}
    \p_v\cap R_\p = \p R_\p.
\end{align*}
Since $\p_v\cap R $ is a prime ideal and therefore it is maximal. $\p_v\cap R_\p$ is the prime ideal of $R_\p$ lying over $\p_v\cap R$, then we have $\p_v\cap R_\p$ is also maximal implying what we want.
Therefore 
\begin{align*}
    \bigcap_{v\in\Omega_\p}\p_v \supseteq \p R_{\p}
\end{align*}
For arbitrary nonzero element $x\in\p R_\p$ and $v\in\Omega_\p$, we have $v(x)\neq 0$ which contradicts to (R2) since we suppose $\Omega_\p$ is an infinite set.


(1.c)

Choose $z^{-1}\in\p$ and (1.a) implies that there exists $v$ such that $\p=R\cap \p_v$ and $v(z)<0$. 
It implies that $R_v\supseteq R_\p$. Therefore $\Omega_\p$ is not an emptyset.

First I claim that 
\begin{align*}
    R_\p = \bigcap_{v\in\Omega_\p} R_v.
\end{align*}
It is easy to see that $R_\p \subseteq \bigcap_{v\in\Omega_\p} R_v$ by definition. Conversely, suppose that $x\in \bigcap_{v\in\Omega_\p} R_v$. Then we obtain that $v'(x)\ge 0$ for arbitrary $v'\in\Omega_\p$.
For arbitrary $v\in (\Omega\setminus \Omega_\p)$, we have $\p\neq R\cap \p_v$. If not, $v$ must contained in $\Omega_\p$. Therefore, by (1.a), we find that $R\cap x^{-1}R$ is not contained in $\p$.
Suppose that $y\in (R\cap x^{-1}R)-\p$, then we find for arbitrary $v\in\Omega$, $v(y),v(xy)\ge 0$ and for arbitrary $v'\in\Omega_\p$, we have $v'(y)=0$. It means that $xy\in R\subseteq R_\p$ and $y\in R-\p$.
Since $v'(R-\p)= v'((R_{v'}-\p_{v'})\cap R)=0$ and $v'(\p)=v'(\p_{v'}\cap R)\ge 1$, then it must have $y\in R-\p$, therefore $x\in R_\p$ meaning that $R_\p \supseteq \bigcap_{v\in\Omega_\p} R_v$.

Select an arbitrary valuation $v\in \Omega_\p$, I claim that $R_v=R_\p$. From the discussion of (1.b), we obtain that 
\begin{align*}
    v_\p\cap R_\p = \p R_p.
\end{align*}
It means that $v(R_\p-\p R_p)= 0$. Hence, we have 
\begin{align*}
    v^2_\p\cap R_\p = \p^2 R_p.
\end{align*}
We notice that $v^2_\p\cap R_\p\supseteq \p^2 R_p$ since $v(\p^2 R_\p)\ge v(\p^2)\ge 2$. Therefore, $v^2_\p\cap R_\p$ must equals $\p^2 R_\p$ or $\p R_\p$. If it equals $\p R_\p$,
then for arbitrary $x\in K^*$, if $v(x)=1$ then we have $x\notin R_\p$, which is wrong. To see that, we choose an arbitrary element $x\in K^*$ such that $v(x)=1$.
Suppose that $\Omega_\p=\{v,v_1,\dots,v_k\}$.
Then we use the weak approximation theorem, there exists an element $y\in K^*$ such that $v(y)=0$ and $v_i(y)=-v_i(x)$.
Then we find that $v(xy)=1$ and $v_i(xy)=0$ meaning that $xy\in R_\p$. Therefore, it must have $v^2_\p\cap R_\p=\p^2 R_\p$. Use same method, we can prove that $v^j_\p\cap R_\p=\p^j R_\p$ for any positive integer $j$. 
It meaning that $R_v= R_\p$. Since select arbitrarily, the uniqueness is obtained.

(2) Let $\p=\p_v\cap R$. It is a prime ideal and hence is maximal ideal. We find that $v(R-\p)=0$ and therefore $v(R_\p)\ge 0$ implying $R_p\subseteq R_v$. 
By (1), we have $R_\p=R_v$. If there exists a different maximal ideal $\mathfrak{q}$ such that $R_v=R_\mathfrak{q}$. Then it means that $\p_v$ is lying over $\mathfrak{q}$, then it must equal to $\p$.

(3) We just need to prove that $R_\p\subseteq R_v$ according to (1). It is quite clear from the construction of (2).

(4.a) We notice that $v_{\p}(I)\neq 0$ if and only if $I\subseteq \p$. If set $\{\mathfrak{p} \in \text{Spm}(R) \mid v_{\mathfrak{p}}(I) \neq 0\}$ is not finite, It means that 
for arbitrary element $x\in I$, there exists infinite valuations $v$ such that $v(x)\neq 0$.
Then by (R2) we must have $I=(0)$ 
contradicting to our condition.

(4.b) Let $J$ denote $\prod_{i=1}^n \mathfrak{p}^{r_{\mathfrak{p_i}}}$ such that for $\p\notin\{\p_1,\dots,\p_n\}$ we have $v_\p(I)=0$ and
$r_\p=v_\p(I)$.
We notice that $I_\p = \p^{r_\p}R_\p$ since
\begin{align*}
    I_\p \subseteq \p^{r_\p}R_\p \text{ and } I_\p \subsetneq \p^{r_\p+1}R_\p
\end{align*}
Then we have 
\begin{align*}
    I= \bigcap_{\p \in \text{Spm}(R)} I_\p =\bigcap_{i=1}^{n} I_{\p_i}= J.
\end{align*}
The last step is because $\p_i$ are pairwise coprime.
(4.c) It is directly deduced by (4.b).

(5) Noetherian is due to (4.c). For intergrality, we consider ideal $(x)$ generated by an arbitrary element $x\in R$.
\end{solution}


\section*{Homework 3}
\subsection*{Exercise 3.1}
Let $L$ be an $R$-lattice contained in the $K$-vector space $V$.
Prove that $L$ is an $R$-lattice supported in $V$ if and only if $FL$ has a $K$-free basis.
\begin{solution}
    Suppose that $L$ is supported in $V$. It means that $\dim_F(FL)=\dim_K(KL)$.  Choose a basis $\{e_i\}$ of $FL$.
    Since $\{e_i\}$ spans $KL$, then if $\{e_i\}$ is not a basis of $KL$, we then have $\dim_K(KL)<\dim_F(FL)$. This is a contradiction.
    Conversely, it is trivial.
\end{solution}


\section*{Homework 4}

\subsection*{Exercise 4.1}
Let $\mathfrak{p}_1, \cdots, \mathfrak{p}_m$ be prime ideals of a ring $R$. Prove that for any ideal $I \subseteq R$, $I \subseteq \bigcup_{i=1}^m \mathfrak{p}_i$ if and only if $I \subseteq \mathfrak{p}_i$ for some $1 \leq i \leq m$.
\begin{solution}
    Suppose that $I \subseteq \bigcup_{i=1}^m \mathfrak{p}_i$. 
    We use induction on $n$ and starting by $n=2$.
    Suppose that $a_1\in \p_1\cap \p_2^c\cap I$ and $a_2\in \p_1^c\cap\p_2\cap I$ where both $a_1$ and $a_2$ are not equal to zero.
    Since $I$ is an ideal, we then have $a_1+a_2\in I$. Clearly, $a_1+a_2$ must not contained in $\p_1$ and $\p_2$. Otherwise, we then have $a_1\in\p_2$ or $a_2\in\p_1$. It is a contradiction.

    Suppose our claim holds for $n-1$. Suppose that there exist an element $a_i\in I\cap \p_i$ such that $a_i$ is not contained in any $I\cap\p_j$ for $j\neq i$. Otherwise, we can reduce it to $n-1$ case.
    Then we consider the elements $a=\sum_{i}\prod_{k\neq i} a_k$. Since $I$ is an ideal then we have $a\in I$. Suppose that $a\in\p_j\cap I$.
    Then we obtain that $\prod_{k\neq j} a_k\in \p_j$. Since $\p_j$ is a prime ideal and $a_i\notin \p_j$ for $i\neq j$, we then have $\prod_{k\neq j} a_k\notin \p_j$.
    It is a contradiction.
    Therefore, our claim holds.
\end{solution}



\subsection*{Exercise 4.2}
Let $R$ be a Dedekind domain, and let $S \subseteq R \setminus \{0\}$ be a multiplicative subset.
\begin{enumerate}
    \item Prove that the ring of fractions $S^{-1}R$ is a Dedekind domain.
    \item Suppose that $S = R \setminus \bigcup_{i=1}^m \mathfrak{p}_i$ for finitely many maximal ideals $\mathfrak{p}_1, \cdots, \mathfrak{p}_m$ of $R$. Prove that $S^{-1}R$ is a PID.
\end{enumerate}
\begin{solution}
    (1) For each $\p\cap S=\emptyset$, we have 
    \begin{align*}
        (S^{-1}R)_{S^{-1}\p} =  S^{-1} R\otimes_R (R\setminus \p)^{-1}= S^{-1}R_\p = R_\p
    \end{align*}
    Since $R$ is Dedekind domain, we then have $R_\p$ is DVR. Therefore, $S^{-1}R$ is also a Dedekind domain.

    (2)
    Since Dedekind domain $S^{-1}R$ only have finite prime ideal $S^{-1}\p_i$. According to Excercise 1.2, it is a PID.
\end{solution}




\subsection*{Exercise 4.3}
Let $R$ be a Dedekind domain, and let $M \neq 0$ be a finitely generated $R$-module. Let $M_{\mathrm{tors}}$ be the torsion submodule of $M$. The invariants factors of $M$ are defined as those of $M_{\mathrm{tors}}$ (cf. Thm. 1.2.27).
\begin{enumerate}
    \item Prove that there is a torsion-free submodule $N \subseteq M$ such that $M = M_{\mathrm{tors}} \oplus N$, and that the isomorphism class of $N$ is uniquely determined by $M$.
    
    We define the \textbf{dimension} of $M$ and the \textbf{Steinitz class} of $M$ as the dimension of $N$ and the Steinitz class of $N$ respectively.
    
    \item Prove that the isomorphism class of $M$ is determined by its invariant factors, dimension and Steinitz class.
\end{enumerate}
\begin{solution}
    (1) Because $M/M_{\text{tros}}$ is a $R$-lattice. It is projective since $R$ is Dedekind domain.
    We have the exact sequence 
    \begin{align*}
        0\to M_{\text{tros}} \to M \to M_{\text{tros}}\to 0
    \end{align*} 
    split. We then have $M = M/M_{\text{tros}}\oplus M_{\text{tros}}$.

    (2) Let $M\simeq M'$. We have $M_{\text{tros}}\simeq M'_{\text{tros}}$ and $M/M_{\text{tros}}\simeq  M'/M'_{\text{tros}}$.
    The torsion-free part is uniquely determined by its Steinitz class and dimension.
    According to our textbook, the invariant factors are uniquely determined by torsion part $M_{\text{tros}}$.
    We just need to prove that for integral ideal $\mathfrak{a}_i$ and $\mathfrak{b}_j$ such that $\mathfrak{a}_i\supseteq \mathfrak{a}_{i+1}$ and $\mathfrak{b}_j\supseteq \mathfrak{b}_{j+1}$ 
    if 
    \begin{align*}
        \bigoplus_{i=1}^n R/\mathfrak{a}_i \simeq \bigoplus_{j=1}^m R/\mathfrak{b}_j
    \end{align*}
    we then have $\mathfrak{a}_i = \mathfrak{b}_i$.
    It is also trivial, since for every prime ideal $\mathfrak{q}$ of $R$, this two module must agree on those localization.

    If $\mathfrak{q}\supseteq \mathfrak{a}_n$ but $\mathfrak{q}\not\supseteq \mathfrak{b}_m$, then we have 
    \begin{align*}
        \left(\bigoplus_{j=1}^m R/\mathfrak{b}_j\right)_\mathfrak{q}
    \end{align*}
    is a zero $R_\mathfrak{q}$-module.
    However, 
    \begin{align*}
       \left(\bigoplus_{i=1}^n R/\mathfrak{a}_i\right)_\mathfrak{q}
    \end{align*}
    is a torsion $R_\mathfrak{q}$-module. It is impossible.
    Therefore, we must have $\mathfrak{a}_n$ and $\mathfrak{b}_m$ shares the same prime factors.
    Moreover, since $R_\mathfrak{q}$ is a DVR, and thus is a PID. According to the classification of modules over PID and 
    the fact that the prime elements on DVR are exactly uniformizers of $R_\mathfrak{q}$. It implies that 
    \begin{align*}
        (\mathfrak{a}_i)_\mathfrak{q}  = (\mathfrak{b}_i)_\mathfrak{q}.
    \end{align*}  
    and $n=m$.
    Our claim holds.
\end{solution}



\subsection*{Exercise 4.4}
Let $R$ be an integral domain with fraction field $F$. Let $L$ be a nonzero $R$-lattice in a finite dimensional $F$-vector space $V$.
Prove that $\dim L$ is equal to the largest positive integer $n$ such that $L$ contains a linearly independent system of $n$ vectors of $V$.

\begin{solution}
    Suppose that $\omega_1,\dots,\omega_n$ is such a system. Suppose that $\omega\in FL$
    can not be linearly represented by $\omega_i$ over $V$. Therefore, $\{\omega,\omega_1,\dots,\omega_n\}$ is linearly independent over $F$.
    Since $R$ is integral domain and $F$ is the fraction field of $R$, there exist an element $r\in R$ such that $\{r\omega,r\omega_1,\dots,r\omega_n\}$
    is included in $L$ and for each of those elements is not zero.
    Then we find a linearly independent system larger then previous one, which is a contradiction.
\end{solution}

\subsection*{Exercise 4.5}
Let $R$ be a Dedekind integral domain with fraction field $F$. Let $L$ be a nonzero $R$-lattice in a finite dimensional $F$-vector space $V$.

Prove that for any nonzero $x \in L$ the following are equivalent:
\begin{enumerate}
    \item $x$ is a primitive vector in $L$ (in the sense of Definition 1.2.17).
    \item the quotient $R$-module $L/Rx$ is torsion-free.
    \item $x$ is a maximal vector in $L$ (in the sense of (1.2.31)).
\end{enumerate}
\begin{solution}
    We first show that (1) implies (2). $x$ is primitive means the exact sequence 
    \begin{align*}
        0\to Rx \to L \to L/Rx\to 0 
    \end{align*}
    splits and $L\simeq Rx\oplus L/Rx$. If $L/Rx$ is not torsion-free,
    then there exists $r\in R$ and $y\in L/Rx$ such that $ry\in Rx$, the direct sum implies that $ry=0$. 
    then $L$ has a torsion element. Consequently, it cannot be a $R$-lattice.

    (2) also implies (1). Since $L/Rx$ as a $R$-lattice, it is projective $R$-module, and thus the exact sequence splits.

    Next, we show that (2) implies (3). Let $\mathfrak{a}$ be the coefficient ideal of $x$. Then we have $\mathfrak{a}x/Rx$ is torsion-free module.
    It also say that 
    \begin{align*}
        \mathfrak{a}/R \simeq R/\mathfrak{a}^{-1}
    \end{align*}
    is a torsion-free module. Actually, we choose an arbitrary element $d$ in determinant ideal $\det(\mathfrak{a})$, we then have $d\mathfrak{a}\subseteq R$.
    It holds if and only if $\mathfrak{a}=R$.
    Therefore, $x$ is maximal.

    We check (3) implies (2). Let $\mathfrak{a}$ be the coefficient ideal of $x$. Suppose that $L/\mathfrak{a}x$ has torsion element $y+\mathfrak{a}x$.
    It means that there exists $r$ such that $ry \in \mathfrak{a}x$ but $y\notin \mathfrak{a}x$.
    We can suppose that $ry = ax$ and $a/r x\notin \mathfrak{a}x$ means $a/r\notin \mathfrak{a}$. Then we have $y = a/r x\notin L$.
    It is a contradiction.

    Therefore, statements are all equivalent.
\end{solution}

\section*{Homework 5}

\subsection*{Exercise 5.1}
Let $I, J$ be fractional ideals of a Dedekind domain $R$. Prove that $I^{-1} \cap J^{-1} = (I + J)^{-1}$.
\begin{solution}
    Let $I^{-1} = \p_1^{-a_1}\cdots \p_n^{-a_n}$ and $J^{-1}=\p_1^{-b_1}\cdots \p_n^{-b_n}$.
    Then $I+J = \p_1^{c_1}\cdots \p_n^{c_n}$ where $c_i=\min\{a_i,b_i\}$.
    Let $v_i$ be the valuations correspond to $\p_i$ respectively.
    Then we have 
    \begin{align*}
        I^{-1}\cap J^{-1} =& \{r\in F\mid v_i(r)\ge \max\{-a_i,-b_i\}\text{ and $w(r)\ge 0$ for others valuation $w$}\}\\
        =& \{r\in F\mid v_i(r)\ge -\min\{-a_i,-b_i\}\text{ and $w(r)\ge 0$ for others valuation $w$}\} \\
        =& (I+J)^{-1}.
    \end{align*}
\end{solution}



\subsection*{Exercise 5.2}
Let $J$ be an integral ideal of a Dedekind domain $R$. Prove that for every fractional ideal $I$ of $R$, one has $I/IJ \cong R/J$ as $R$-modules.
\begin{solution}
    Let $\p$ be an arbitrary prime ideal. Then we have isomorphim of $R_\p/J_\p$-module
    \begin{align*}
        (I/IJ)_\p = I_\p/(I_\p J_\p) \simeq R_\p/J_\p.
    \end{align*}
    The isomorphism is given by $R_\p\simeq I_\p$. 
    Interpret $I/IJ$ and $R/J$ be the $R/J$-module. 
    According to Structure Theorem for Artinian Modules, we have 
    \begin{align*}
        I/IJ \cong R/J.
    \end{align*}
    Hence it is an isomorphism as $R$-modules.
\end{solution}


\subsection*{Exercise 5.3}
Prove that the existence statement in Thm.1.2.27 can be deduced from Thm.1.2.39.
\begin{solution}
    Let $\{x_1,\cdots,x_n\}$ be the generators of torsion $R$-module $M$. We consider the homomorphism 
    \begin{align*}
        R^{n} &\to M \\
        (r_i)_i&\mapsto \sum r_ix_i
    \end{align*}
    which is a epimorphism. Suppose its kernel is $N$. Clearly, $N$ is a torsion-free module with finite generators since $R$ is Noetherian. Therefore $N$ is a projective $R$-module.
    Consider the short exact sequence 
    \begin{align*}
        0\to N \to R^n \to M \to 0
    \end{align*}
    It implies that 
    \begin{align*}
        K^n \simeq N \otimes_R K.
    \end{align*}
    Since functor $\otimes_R K$ is exact and $M\otimes_R K$ is zero.

    Then there is a basis such that 
    \begin{align*}
        R^n = \a_1y_1 +\a_2y_2 +\cdots+ \a_ny_n\\
        N = \a_1\b_1y_1+\a_2\b_2y_2+\cdots + \a_n\b_ny_n
    \end{align*}
    and $\b_i\supseteq \b_{i+1}$.
    Then we obtain that 
    \begin{align*}
        R^n/N \simeq \bigoplus_i \a_i/(\a_i\b_i) \simeq  \bigoplus_i R/\b_i.
    \end{align*}
\end{solution}

\subsection*{Exercise 5.4}
 Prove the following facts about metabolic bilinear spaces over a field $F$:
\begin{enumerate}
    \item The hyperbolic spaces $m.\mathbb{H}$ are metabolic.
    \item For any symmetric bilinear spaces $V, V'$, if two of the three spaces $V, V'$ and $V \perp V'$ are metabolic, then so is the third. In particular, an orthogonal sum of metabolic spaces is metabolic.
    \item For any nonsingular symmetric bilinear space $(V, b)$, the space $(V, b) \perp (V, -b)$ is metabolic.
    \item Let $V$ be a 2-dimensional nonsingular symmetric space. Then $V$ is metabolic if and only if it is isotropic.
    \item Every 2-dimensional metabolic space is isomorphic to $\mathbb{H}$ or $\langle a, -a \rangle_{\mathrm{bil}}$ for some $a \neq 0$.
    \item If $\mathrm{char}(F) \neq 2$, then every metabolic bilinear space is hyperbolic.
\end{enumerate}
\begin{solution}
    (1) Clearly, $m.\mathbb{H}$ it is not singular of even dimension. Let $\{x_i,y_i\}$ be the basis of $i^{\text{th}}$ hyperbolic subspace of $m.\mathbb{H}$,
    then we let $V=\text{span}(x_i)$. It is easy to check that $V$ is totally isotropic and $\dim(V)=m$.

    (2) Suppose that $V$ and $V'$ are metabolic. Clearly, $V\perp V'$ is not singular of even dimension. Let $W$ and $W'$ be the totally isotropic subspace of $V$ and $V'$ respectively and 
    hance have half dimension. Then $W\perp W'$ also be the totally isotropic subspace of $V\perp V'$ with a half of dimension. Therefore, $V\perp V'$ is totally isotropic.

    Suppose that $V$ and $V\perp V'$ is metabolic. Then we have $\dim(V')= \dim(V\perp V')-\dim(V)$ is even.
    And $V'$ must be not singular. Otherwise, $V\perp V'$ will be singular. Let $W$ be the totally isotropic subspace of $V\perp V'$ with a half dimension.
    Then we have $W = (W\cap V)\perp (W\cap V')$. I claim that $\dim(W\cap V)\le \dim(V)/2$. If $\dim(W\cap V)> \dim(V)/2$, then we choose a basis of $W\cap V$ and extend it to be a basis of $V$.
    Let we denote it by $\{e_1,\dots,e_n,\dots,e_{2n}\}$. Write the Gram matrix over this basis and using Laplace expension, we deduce that $V$ is singular.
    It is a contradiction. Therefore, we have $\dim(V'\cap W)\ge \dim(W)-\dim(V)/2 = \dim(V')/2$.
    
    (3) Let $e_i$ be a standard basis of $(V,b)$. Then it is also a standard basis of $(V,-b)$. 
    Consider the basis $(e_i,e_i),(e_i,-e_i)$ of $V\oplus V$. Check subspace generated by $(e_i,e_i)$ is totally isomorphic subspace with half dimension of $(V,b)\perp (V,-b)$.

    (4) If it is isotropic, then it has Gram matrix $$ \begin{pmatrix} 0 & a \\ a & c \end{pmatrix} $$ with $a\in F^*$. 
    Consider base change matrix 
    
    It is a hyperbolic and thus is a metabolic.
    Conversely, if it is a metabolic, it has the Gram matrix $$ \begin{pmatrix} 0 & a \\ a & c \end{pmatrix} $$
    It implies that $V$ is isotropic.

    (5) If $c=0$, then $V$ is isomorphic to $\mathbb{H}$. If $c\neq 0$, let $\{e_1,e_2\}$ be the basis corresponding to this matrix.
    Then we consider the basis $\{e_2,e_2-a^{-1}ce_1\}$. It has the Gram matrix
    $$ \begin{pmatrix} c & 0 \\ 0 & -c \end{pmatrix} $$
    The statements holds.

    (6)  Let $U$ be the totally isotropic subspace of $V$ with half dimension.
    According to our textbook, metabolic space can be decomposit to $V_1\perp \dots \perp V_n$ where $V_i=Fx_i\otimes Fy_i$ is a subspace of dimension 2 where $(x_i,y_i)=1$ and $x_i\in U$.
    Therefore, $V_i$ is isotropic. 
    Suppose that $V$ is isotropic to the second case of (5). 
    We consider basis $\{x_i+y_i,x_i-y_i\}$. It has the Gram matrix
    $$ \begin{pmatrix} 0 & 2c \\ 2c & 0 \end{pmatrix} $$
    Since $\text{char}(F)\neq 2$, then we have $V_i$ is hyperbolic space and $V\simeq n.\mathbb{H}$.
\end{solution}


\section*{Homework 6}

In the following exercises, $F$ denotes a field of characteristic $\text{char}(F)\neq 2$. 
In principle, the following exercises can be done by using only results stated in \S 2.1.

\subsection*{Exercise 6.1}
Let $V$ be a nonsingular quadratic space over $F$. Prove that $V$ is isotropic if and only if $V$ represents $\mathbb{H}$.
\begin{solution}
    Suppose that $Q(x)=0$. Since $V$ is not singular, then there exists $y\in V$ such that $V = V'\perp(Fx\oplus Fy)$. Since characteristic is not equal 2, then we have isotropic subspace $Fx\oplus Fy$ is hyperbolic.
    Therefore we have representation $\varphi:\mathbb{H} \to Fx\oplus Fy$.
    Converse is trivial.
\end{solution}


\subsection*{Exercise 6.2}
Let $V$ be a quadratic space over $F$ and $\alpha \in F^\times$. Prove that the following statements are equivalent:
\begin{enumerate}
    \item[(i)] $V$ represents the 1-dimensional space $\langle \alpha \rangle$.
    \item[(ii)] There is a splitting $V = \langle \alpha \rangle \perp W$.
\end{enumerate}
Prove that if $V$ is nonsingular, then the above conditions are also equivalent to the following:
\begin{enumerate}
    \item[(iii)] The quadratic space $V \perp \langle -\alpha \rangle$ is isotropic.
\end{enumerate}

\begin{solution}
    We show that (1) implies (2). Since $\alpha\in F^*$, then $\langle \alpha \rangle$ is not singular.
    Then according to the textbook, we have $V\simeq \langle \alpha\rangle \perp W$.
    Converse is trivial.
    
    Suppose that $V$ is nonsingular. Then clearly (ii) implies (iii) by considering the $\langle\alpha\rangle\perp \langle-\alpha\rangle$.
    Conversely, $V\perp \langle -\alpha\rangle$ is also nonsingular. Then according to Excercise 6.1, we have $V\perp \langle -\alpha \rangle$ represents a hyperbolic space $W\simeq\mathbb{H}$ such that $V\perp \langle-\alpha\rangle = W\perp W^{\perp}$.
    Since $\text{char}(F)\neq 2$, $W$ represents $$ \begin{pmatrix} \alpha & 0 \\ 0 & -\alpha \end{pmatrix} $$
    Using Witt cancellation then we have $V\simeq \langle \alpha \rangle\perp W^{\perp}$.
    Therefore, $V$ represents the 1-dimensional space $\langle \alpha \rangle$.
\end{solution}

\subsection*{Exercise 6.3}
Prove that Witt cancellation holds for all quadratic spaces in characteristic $\neq 2$. More precisely, for any (possibly singular) quadratic spaces $V, V_1, V_2$ over $F$, prove that $V \perp V_1 \cong V \perp V_2$ if and only if $V_1 \cong V_2$.
\begin{solution}
The direction ($\Longleftarrow$) is trivial. We prove ($\Longrightarrow$).

Suppose there is an isomorphsim $V \perp V_1 \cong V \perp V_2$.

Since $\mathrm{char}(F) \neq 2$, the relation $B(x,x) = 2Q(x)$ guarantees that the radical of any quadratic space is a totally singular space. This ensures the existence of a valid orthogonal decomposition for any space $U$:
\[ U = \mathrm{rad}(U) \perp U_{nd} \]
where $U_{nd}$ is a strictly nonsingular subspace.

An isometry must preserve the radical. Applying this to our spaces gives:
\[ \mathrm{rad}(V) \perp \mathrm{rad}(V_1) \cong \mathrm{rad}(V) \perp \mathrm{rad}(V_2) \]
For totally singular spaces, isometry is entirely determined by dimension. Taking dimensions on both sides and canceling $\dim(\mathrm{rad}(V))$ yields $\dim(\mathrm{rad}(V_1)) = \dim(\mathrm{rad}(V_2))$. Thus:
\[ \mathrm{rad}(V_1) \cong \mathrm{rad}(V_2) \]

Furthermore, the isometry also preserves the nonsingular orthogonal complements:
\[ V_{nd} \perp (V_1)_{nd} \cong V_{nd} \perp (V_2)_{nd} \]
Because all components here are strictly nonsingular and $\mathrm{char}(F) \neq 2$, we are permitted to invoke the standard Witt's Cancellation Theorem (which relies on non-degenerate reflections) to cancel $V_{nd}$ from both sides. This yields:
\[ (V_1)_{nd} \cong (V_2)_{nd} \]

Finally, we reconstruct $V_1$ and $V_2$ by orthogonally summing their respective isometric parts:
\[ V_1 = \mathrm{rad}(V_1) \perp (V_1)_{nd} \cong \mathrm{rad}(V_2) \perp (V_2)_{nd} = V_2 \]
This completes the proof.
\end{solution}

\section*{Homework 7}

\subsection*{Exercise 7.1}
Let $(V,\langle \cdot,\cdot \rangle)$ be a nonsingular bilinear space and let $x_1,\cdots,x_n$ be a basis of $V$. Let $y_1,\cdots,y_n \in V$ be the dual basis of $x_1,\cdots,x_n$ with respect to the given bilinear form $\langle \cdot,\cdot \rangle$. Prove
\[
    \det(\langle x_i, x_j \rangle) \det(\langle y_i, y_j \rangle) = 1.
\]
\begin{solution}
    Suppose that $(x_i)=A(y_i)$. Then we have $x_i = \sum_{j}a_{ij}y_j$.
    Since $(y_i)$ is the dual basis of $(x_i)$, then we have 
    \begin{align*}
        \langle x_i,x_j\rangle = \langle \sum_{k}a_{ik}y_k,x_j\rangle = a_{ij}.
    \end{align*}
    Next, we consider $A^* (x_i) = \det(A)(y_i)$. We let $A^*=(a^*_{ij})$.
    \begin{align*}
        \det(A)\langle y_i,y_j\rangle = \langle y_i,\sum_{k}a^*_{jk}x_k\rangle = a_{ji}^*.
    \end{align*}
    We find that 
    \begin{align*}
        \det(\langle x_i, x_j \rangle) \det(\langle y_i, y_j \rangle) = \det(A)\det(A^*)\det(A^{-n})=1
    \end{align*}
\end{solution}


\subsection*{Exercise 7.2}
Consider the quadratic $\mathbb{Z}$-lattice $(L, Q)$ with $L = \mathbb{Z}^{\oplus 2}$ and $Q(x,y) = 2xy$.
Prove that $L$ cannot be decomposed as an orthogonal sum of two 1-dimensional lattices.

\begin{solution}
    Let $e_1,e_2$ denotes $(1,0)$ and $(0,1)$ respectively. Then we have 
    \begin{align*}
        \langle ae_1+be_2,c e_1+de_2 \rangle =2(bc+ad).
    \end{align*}
    If $ae_1+be_2,c e_1+de_2$ are orthogonal, then we must have $bc+ad=0$.
    Hence, $ae_1+be_2,c e_1+de_2$ must be the generators of $L$, which means the relative determinant must be $\pm 1$, which is $ad-bc=\pm 1$.
    Then we have 
    \begin{align*}
        2bc = \mp 1.
    \end{align*}
    It is impossible, since we must have $b,c\in\mathbb{Z}$.
\end{solution}

\section*{Homework 9}

\subsection*{Exercise 9.1}
 Let $(V, Q)$ be a quadratic space over a field $F$ of characteristic $\neq 2$.
\begin{enumerate}
    \item Check that for any anisotropic vector $x \in V$, the symmetry $\tau_x$ defined in (2.3.1) is an isometry of $(V, Q)$.
    \item Assume that $(V, Q)$ is nonsingular. Without using the Cartan--Dieudonné theorem, prove that $\det(\tau_x) = -1$ for any anisotropic vector $x \in V$ and that $\det(\sigma) = \pm 1$ for every $\sigma \in \mathrm{O}(V)$.
\end{enumerate}
\begin{solution}
    (1)
    We let $B(x,y)=Q(x+y)-Q(x)-Q(y)$.
    Clearly, $\tau_x$ is invertible since $\tau_x^2 = \text{id}$. 
    On one hand, we have
    \begin{align*}
        B\left(y,\frac{-B(x,y)}{Q(x)}x\right) = Q(\tau_x(y)) - Q(y) - \frac{B^2(x,y)}{Q^2(x)}Q(x).
    \end{align*} 
    On the other hand we have 
    \begin{align*}
        B\left(y,\frac{-B(x,y)}{Q(x)}x\right) = -\frac{B^2(x,y)}{Q(x)}.
    \end{align*}
    Then we obtain that $Q(\tau_x(y))=Q(y)$

    (2) Since $V$ is not singular, there exists an orthogonal decomposition $V = \langle x\rangle+ \langle x\rangle^\perp$. Therefore, the linear transformation $\tau_x$ gives
    matrix of diagonal form $\text{diag}(-1,1,\dots,1)$. Then we obtain that $\det(\tau_x)=-1$.

    To prove that $\det(\sigma)=\pm 1$. Take $\{e_i\}$ be the orthogonal basis whose Gram matrix is given by $G$. Let $M$ be the matrix representing the linear transformation $\sigma$.
    Then we have 
    \begin{align*}
        M^TGM = G
    \end{align*}
    and hence $\det(G)\det^2(M)=\det(G)$. We then find that $\det(M)=\pm 1$.
\end{solution}

\subsection*{Exercise 9.2}
 Let $M$ be a subgroup of the additive group $\mathbb{R}$ endowed with the usual Euclidean topology. Prove that the following assertions are equivalent:
\begin{itemize}
    \item[(i)] $M$ is not dense in $\mathbb{R}$.
    \item[(ii)] There exists a real number $\varepsilon > 0$ such that $M \cap \left]-\varepsilon, \varepsilon\right[ = \{0\}$.\\
    (Here $\left]-\varepsilon, \varepsilon\right[$ denotes the open interval $\{x \in \mathbb{R} \mid -\varepsilon < x < \varepsilon\}$.)
    \item[(iii)] $M$ is discrete in $\mathbb{R}$, i.e., the subspace topology on $M$ is the discrete topology.
    \item[(iv)] $M = r\mathbb{Z}$ for some $r \in \mathbb{R}$.
\end{itemize}
\begin{solution}
    We first show that (1) implies (2). Suppose that $M$ is dense on $0$. Then for arbitrary $x\in M$ and $y\in \mathbb{R}$, there exists $\varepsilon\in M$ such that $|x-y|<\varepsilon$.
    If $x>y$, it implies that $x>x-\varepsilon>y$ and if $x<y$, it implies that $x<x+\varepsilon<y$. In any case, we find an element $x\pm\varepsilon$ is closer than $x$ to $y$.
    Equivalently we write $d(M,y)=0$. It contradicts to (1).

    Clearly, (2) also implies (1) since $M$ is not dense on zero.

    Hence, (2) implies (3). Because $M$ must be nowhere dense. If $M$ is dense on $x$, for arbitrary $\varepsilon>0$ such that sufficiently small, then there exists an $x+\eta_1,x+\eta_2\in M$ such that $|\eta_1-\eta_2|<\varepsilon$.
    It follows that $M$ is dense on zero, which contradicts to (2).

    (3) implies (2) also holds by defenition.

    (4) implies (2) is trivial.

    Finally, we check that (3) implies (4). We consider the set of distence between elements in $M$, denoted by 
    \begin{align*}
        D=\{d(x,y)\mid x,y\in M\}.
    \end{align*}
    Since $M$ is discrete, we can find a miminal value $\delta$ of $D$. Clearly, we have $\delta\mathbb Z\subseteq M$. 
    It there exists an element $y\in M$ is not contained in $\delta \mathbb{Z}$, then we find that $d(y,\delta \mathbb{Z})<\delta$. It contradicts to the minimality of $\delta$.
    Therefore, we then have $M=\delta \mathbb{Z}$.
\end{solution}


\subsection*{Exercise 9.3}
 Let $K$ be a field and let $v$ be an exponential valuation on $K$ with valuation ring $A$. Let $(\hat{K}, \hat{v})$ be the completion of the valued field $(K, v)$.
\begin{enumerate}
    \item Let $\hat{A}$ denote the valuation ring of $\hat{v}$ in $\hat{K}$. Prove that $\hat{A}$ equals the topological closure of $A$ in $\hat{K}$.
    \item Show that $\hat{v}(\hat{K}^\times) = v(K^\times)$ and that the natural map $A \to \hat{A}$ induces an isomorphism on their residue fields. In particular, if $v$ is discrete (resp. discrete and normalized), then so is $\hat{v}$.
\end{enumerate}
\begin{solution}
    \textbf{Maybe we should show stability first.}

    (1) Let $\bar{A}$ denote the topological closure of $A$. Let $||$ denote the absolute valuate determined by $v$.
    Let $(x_i)$ be a Cauchy sequence over $K$. Suppose that $(x_i)\in \hat{A}$, it means that $\lim_{n\to\infty}|x_n|\le 1$. 
    Then it is to said that there exists an integer $N>0$ for any $n>N$ we have $x_n\in A$ since stability.  The we consider constant sequences $(x_{N+k})_i$.
    For arbitrary $\varepsilon>0$, there exists an integer $N$ such that for $n>N$ we have 
    \begin{align*}
        |x_n-x_{N+k}|<\varepsilon
    \end{align*}
    for arbitrary $k>0$. It equivalently say given any $\varepsilon>0$, we find an element $(x_{N+k})_i\in A$ such that $\|(x_{N+k})_i - (x_i)_i\|<\varepsilon$ where $\|\|$ denotes the absolute valuate induced by $\hat{v}$.
    Thus we have $\hat{A}\subseteq \bar{A}$.

    Conversely, since $\hat{A}$ is closed in $\hat{K}$ and it contains $A$, we then have $\bar{A}\subseteq \hat{A}$.
    Therefore, we obtain that $\bar{A}=  \hat{A}$

    (2) We just need to show that any Cauchy sequence $(x_i)$ admits a stable sequence of valuation $v(x_i)$ if $(x_i)$ is not the zero over $\hat{K}$.
    To see that we use the strongly triangular inequality. Let $a = \lim_{n\to \infty}|x_n|$. There exists $n>0$ such that for any $k$ we have 
    \begin{align*}
        |x_n-x_{n+k}|<a/2.
    \end{align*}
    We can also choose sufficiently large of $n$ such that
    \begin{align*}
        ||x_{n+k}|-a|<a/2
    \end{align*} 
    In this sense we have $|x_{n+k}|>a/2>|x_n-x_{n+k}|$, then 
    \begin{align*}
        |x_{n}| = |x_{n}-x_{n+k}+x_{n+k}| = |x_{n+k}|.
    \end{align*}
    Therefore we show that $|x_n|$ is stable and thus $\hat{v}(\hat{K}^\times) = v(K^\times)$. 

    Let $\mathfrak{m}$ and $\hat{\mathfrak{m}}$ denote the maximal prime ideal of $A$ and $\hat{A}$ respectively. 
    Consider the homomorphism 
    \begin{align*}
        A \to \hat{A}/\hat{\mathfrak{m}}.
    \end{align*}
    It is epimorphic. For arbitrary Cauchy sequence $(x_i)$ such that $\lim_{n\to\infty}|x_n| = 1$, we suppose it is stable at $k^{\text{th}}$.
    Hence we can suppose that $|x_{k+n}-x_k|<1/2$, then we find that $(x_n-x_k)_n\in \hat{\mathfrak{m}}$ and therefore $(x_n)$ have preimage $(x_k)_n$.
    Its kernal is $A\cap \hat{\mathfrak{m}} = \mathfrak{m}$.

    Our claim holds.
\end{solution}



\subsection*{Exercise 9.4}
Let $R$ be a proper subring of a field $K$ such that for every $x \in K^\times$, either $x \in R$ or $x^{-1} \in R$. (We say that $R$ is a \textbf{valuation ring} in $K$.)
\begin{enumerate}
    \item Prove that $R$ is integrally closed.
    \item Let $v$ be a nontrivial exponential valuation on a field $K$. Show that the ring $\mathcal{O}_v = \{x \in K \mid v(x) \ge 0\}$ is a valuation ring in the above sense, i.e., for all $x \in K^\times$, either $x \in \mathcal{O}_v$ or $x^{-1} \in \mathcal{O}_v$. Deduce that $\mathcal{O}_v$ is integrally closed.
\end{enumerate}
\begin{solution}
    (1) Suppose that it is not integrally closed. Equivalently, there exits an integral element $x\in K\setminus R$. Suppose that 
    \begin{align*}
        x^n + a_{n-1}x^{n-1}+\dots+a_0 = 0
    \end{align*}
    Then we have 
    \begin{align*}
        a_0x^{-n} + a_1x^{-n+1}+\dots+1 = 0
    \end{align*}
    Hence we have 
    \begin{align*}
        x = -(a_0x^{-n+1}+a_1x^{-n+2}+\dots+a_{n-1}).
    \end{align*}
    It implies that $x\in R$.

    (2) If $x\notin \o_v$, then we have $v(x)<0$. Therefore, we must have $v(x^{-1})>0$ and $x^{-1}\in\o_v$. The intergrality of $\o_v$ follows.
\end{solution}


\subsection*{Exercise 9.5}
Let $|\cdot|_1, \cdots, |\cdot|_n$ be pairwise inequivalent non-archimedean absolute values of a field $K$. Prove that for all $a_1, \cdots, a_n \in K^\times$ there exists an element $x \in K^\times$ such that $|x|_i = |a_i|_i$ for each $i = 1, \cdots, n$.
\begin{solution}
    We use Weak Approximation Theorem. Let $\varepsilon<\min_i\{|a_i|_i\}$, then there exist an element $x\in K$ such that 
    \begin{align*}
        |x-a_i|_i<\varepsilon.
    \end{align*}
    Next, we consider 
    \begin{align*}
        |x|_i = |x-a_i+a_i|_i = |a_i|_i
    \end{align*}
    since $|x-a_i|_i<|a_i|_i$.
\end{solution}



\subsection*{Exercise 9.6}
Let $(K, |\cdot|)$ be a complete non-archimedean valued field. Let $(a_n)_{n \in \mathbb{N}}$ be a sequence in $K$. Prove that the series $\sum_{n=0}^\infty a_n$ converges in $K$ if and only if $\lim_{n \to \infty} a_n = 0$ in $K$.
\begin{solution}
    Suppose that $\sum a_n$ converges. Then $|a_{n+k+1}-a_{n+k}|<\varepsilon$ therefore $|a_{n+k}-a_n|=|a_{n+k}-a_{n+k-1}+\dots+a_{n+1}-a_n|<\varepsilon$ which implies that $a_n$ converges.

    Conversely, $|a_n+\dots+a_{n+k}|\le \max_{k}\{|a_{n+k}|\}<\varepsilon$. $\sum a_n$ converges.
\end{solution}


\subsection*{Exercise 9.7}
Let $(K, |\cdot|)$ be a henselian (non-archimedean valued) field.
\begin{enumerate}
    \item Let $f(T) = a_2T^2 + a_1T + a_0$ be a polynomial of degree $2$ over $K$.\\
    Prove that if $f$ is irreducible over $K$, then $|a_1|^2 \le |a_0| \cdot |a_2|$.
    \item Let $(V, Q)$ be an anisotropic quadratic space over $K$. Prove that
    \[ |\langle x, y \rangle_Q|^2 \le |Q(x)| \cdot |Q(y)| \quad \text{for all } x, y \in V. \]
    (Here $\langle x, y \rangle_Q := Q(x+y) - Q(x) - Q(y)$.)
\end{enumerate}
\begin{solution}
    (1) Suppose that $|a_1|^2>|a_0a_2|$, it equivalently say that $a_0a_1^{-2}a_2\in \p$ where $R$ is the valuation ring induced by absolute value.
    Without loss of generality, let $f(T)$ be the primitive polynomial. 
    Let $T = a_0a_1^{-1}X$, we then obtain that 
    $$f(X) = a_0 \left( \frac{a_2 a_0}{a_1^2} X^2 + X + 1 \right)$$
    Let $$g(X) =\frac{a_2 a_0}{a_1^2} X^2 + X + 1 $$
    which is the polynomial with coefficients in $R$. Consider that
    \begin{align*}
        \bar{g}(X) = X+1 \bmod \p
    \end{align*}
    Then $\bar{g}(X)$ has a simple root since $\bar{g}'(X)= 1$. Then by Hensel's Lemma, it can be lifted to $R$.
    However, $g(X)$ is irreducible and thus it can not have a root.

    (2) Let $a_2= Q(x),a_0=Q(y)$ and $a_1= \langle x, y \rangle_Q$. 
    Then $Q(Tx+y) = a_2T^2+a_1T+a_0$. Since $V$ is anisotropic, then $Q(Tx+y)$ is irreducible over $K$.
    Then according to (1) is what we want.
\end{solution}


\section*{Homework 10}

\subsection*{Exercise 10.1}
Let $K = \mathbb{R}$ or $K = \mathbb{C}$, and let $|\cdot|$ be the usual absolute value on $K$. Let $s$ be a positive real number. Prove that $|\cdot|^s$ is an absolute value on $K$ if and only if $0 < s \leq 1$.

\begin{solution}
    AV1 and AV2 is quite clear. We just show AV3.
    Suppose that $|\cdot|^s$ is an absolute value, we consider that 
    \begin{align*}
        2\ge 2^s
    \end{align*}
    It menas that $0<s\le 1$. Suppose that $0<s\le 1$, we just need to prove that 
    \begin{align*}
        1+x^s\ge (1+x)^s
    \end{align*}
    for $0<x< 1$. Consider function $g(x)=1+x^s-(1+x)^s$ with 
    \begin{align*}
        g'(x) = sx^{s-1} - s(1+x)^{s-1}>0
    \end{align*}
    Therefore $g(x)>g(0)=0$. Our claim holds.
\end{solution}

\section*{Homework 11}

\subsection*{Exercise A.11.1}

Let $\mathbb{F}$ be a finite field.

\begin{enumerate}
    \item Prove that every quadratic space of dimension $\geq 3$ over $\mathbb{F}$ is isotropic.

    You should not use Proposition 4.2.7 for this proof.

    \item Prove that $\mathbb{F}$ has separable $2$-dimension $\leq 1$; cf. Definition 4.2.5.
\end{enumerate}

The above statements hold even if $\operatorname{char}(\mathbb{F})=2$. If you like, you can only treat the case $\operatorname{char}(\mathbb{F})\neq 2$.

\begin{solution}
    We assume that $\operatorname{char}(\mathbb{F})\neq 2$ and $q=|\mathbb{F}|$.

    (1)First we prove a simple counting lemma. Let $a,b\in \mathbb{F}^{\times}$. Then the binary quadratic form $ax^2+by^2$ represents every element of $\mathbb{F}$. Indeed, fix $t\in \mathbb{F}$ and consider \[ A=\{ax^2:x\in \mathbb{F}\},\qquad B=\{t-by^2:y\in \mathbb{F}\}. \] Since $\mathbb{F}$ has odd cardinality, the set of squares in $\mathbb{F}$, including $0$, has cardinality $(q+1)/2$. Hence $|A|=|B|=(q+1)/2$. Since $|A|+|B|=q+1>q$, the two subsets $A$ and $B$ must intersect. Thus there exist $x,y\in \mathbb{F}$ such that $ax^2=t-by^2$, or equivalently $ax^2+by^2=t$.
    
    Now let $(V,Q)$ be a quadratic space over $\mathbb{F}$ with $\dim V\geq 3$. Since $\operatorname{char}(\mathbb{F})\neq 2$, the quadratic form $Q$ is diagonalizable. Thus, for a suitable basis, we may write \[ Q\simeq \langle a_1,\dots,a_r,0,\dots,0\rangle \] with $a_i\in \mathbb{F}^{\times}$. If a zero coefficient occurs, then the corresponding basis vector is a nonzero isotropic vector, so $Q$ is isotropic. Hence we may assume that no zero coefficient occurs. Since $\dim V\geq 3$, the restriction of $Q$ to the subspace spanned by the first three basis vectors has the form \[ ax^2+by^2+cz^2 \] with $a,b,c\in \mathbb{F}^{\times}$. By the counting fact above, the binary form $ax^2+by^2$ represents $-c$. Therefore there exist $x,y\in \mathbb{F}$ such that $ax^2+by^2=-c$. It follows that \[ Q(x,y,1,0,\dots,0)=ax^2+by^2+c=0. \] The vector $(x,y,1,0,\dots,0)$ is nonzero, so $Q$ is isotropic. This proves that every quadratic space of dimension at least $3$ over $\mathbb{F}$ is isotropic. 
    
    
    (2)It remains to prove that $\operatorname{sd}_2(\mathbb F)\leq 1$. Let
$L/\mathbb F$ be a finite separable extension, and let $E/L$ be a quadratic
separable extension. Since $\mathbb F$ is finite, $L$ is again a finite field.
Assume $\operatorname{char}(\mathbb F)\neq 2$. Then $E=L(\sqrt d)$ for some
$d\in L^\times\setminus L^{\times 2}$.

We show that $N_{E/L}:E^\times\to L^\times$ is surjective. Let
$a\in L^\times$. Consider the quadratic form
\[
q(x,y,z)=x^2-dy^2-az^2
\]
over $L$. By the first part, this three-dimensional quadratic form is
isotropic. Hence there exists a nonzero vector $(x,y,z)\in L^3$ such that
\[
x^2-dy^2=az^2.
\]
We claim that $z\neq 0$. If $z=0$, then $x^2=dy^2$. Since
$d\notin L^{\times 2}$, this forces $x=y=0$, contradicting that
$(x,y,z)$ is nonzero. Thus $z\neq 0$, and after dividing by $z^2$ we get
\[
\left(\frac{x}{z}\right)^2
-
d\left(\frac{y}{z}\right)^2
=a.
\]
But the left-hand side is
\[
N_{E/L}\left(\frac{x}{z}+\frac{y}{z}\sqrt d\right).
\]
Therefore $a$ lies in the image of $N_{E/L}$. Since $a\in L^\times$ was
arbitrary, the norm map is surjective. Hence, by Definition 4.2.5,
$\operatorname{sd}_2(\mathbb F)\leq 1$.
\end{solution}

\subsection*{Exercise A.11.2}

Let $K$ be a field of characteristic $\neq 2$, and let $(V,\varphi)$, $(U,\psi)$ be nonsingular quadratic spaces of positive dimension over $K$.

\begin{enumerate}
    \item Prove that $(V,\varphi)$ is hyperbolic if and only if $\dim V$ is even and there exists a totally isotropic subspace $W\subseteq V$ such that $\dim V=2\dim W$.

    \item Prove that the quadratic space $(U,\psi)\perp (U,-\psi)$ is hyperbolic.

    \item Prove that $(V,\varphi)\cong (U,\psi)$ if and only if $(V,\varphi)\perp (U,-\psi)$ is hyperbolic.
\end{enumerate}
\begin{solution}
    (1) Let $B(x,y)=\varphi(x+y)-\varphi(x)-\varphi(y)$. Since $(V,\varphi)$ is nonsingular and $\operatorname{char}K\neq 2$, $B$ is nondegenerate. Suppose first that $(V,\varphi)$ is hyperbolic. Then $V=H_1\perp\cdots\perp H_n$, where each $H_i$ is a hyperbolic plane. Choose a hyperbolic basis $e_i,f_i$ of $H_i$. Then $W=\langle e_1,\dots,e_n\rangle$ is totally isotropic, and $\dim V=2n=2\dim W$. 
    
    Conversely, suppose that $W\subseteq V$ is totally isotropic and $\dim V=2\dim W$. Since $\varphi|_W=0$ and $\operatorname{char}K\neq 2$, one has $B|_{W}=0$, hence $W\subseteq W^\perp$. By nondegeneracy, $\dim W^\perp=\dim V-\dim W=\dim W$, so $W=W^\perp$. We prove that $V$ is hyperbolic by induction on $\dim W$. Choose $0\neq e\in W$. Since $B$ is nondegenerate, there exists $y\in V$ with $B(e,y)\neq 0$. For any $a\in K$, $\varphi(y+ae)=\varphi(y)+aB(y,e)$, because $\varphi(e)=0$. Taking $a=-\varphi(y)/B(y,e)$, we get $f=y+ae$ with $\varphi(f)=0$ and $B(e,f)\neq 0$. After scaling $f$, assume $B(e,f)=1$. Thus $H=\langle e,f\rangle$ is a hyperbolic plane. Since $H$ is nonsingular, $V=H\perp H^\perp$. Set $W'=W\cap H^\perp=\{w\in W:B(w,f)=0\}$. The map $W\to K$, $w\mapsto B(w,f)$, is nonzero because $B(e,f)=1$; hence $\dim W'=\dim W-1$. Also $W'$ is totally isotropic and $\dim H^\perp=\dim V-2=2(\dim W-1)=2\dim W'$. By induction, $H^\perp$ is hyperbolic. Therefore $V=H\perp H^\perp$ is hyperbolic.

    (2)Consider $(U,\psi)\perp (U,-\psi)$ on $U\oplus U$. Its quadratic form is \[ q(u_1,u_2)=\psi(u_1)-\psi(u_2). \] Let \[ \Delta=\{(u,u):u\in U\}\subseteq U\oplus U. \] Then $q(u,u)=0$ for all $u\in U$. Moreover, for the associated bilinear form, \[ B_q((u,u),(v,v))=B_\psi(u,v)-B_\psi(u,v)=0, \] so $\Delta$ is totally isotropic. Also \[ \dim(U\oplus U)=2\dim U=2\dim \Delta. \] By part $(1)$, $(U,\psi)\perp (U,-\psi)$ is hyperbolic.

    (3) Assume additionally that $\dim V=\dim U$.

If $(V,\varphi)\cong (U,\psi)$, then
\[
(V,\varphi)\perp (U,-\psi)
\cong
(U,\psi)\perp (U,-\psi),
\]
which is hyperbolic by part $(2)$.

Conversely, suppose that $(V,\varphi)\perp (U,-\psi)$ is hyperbolic. Then
$(V,\varphi)$ and $(U,\psi)$ have the same Witt class. By the Witt
decomposition theorem, their anisotropic parts are isomorphic. Hence we may
write
\[
(V,\varphi)\cong H^r\perp A,\qquad
(U,\psi)\cong H^s\perp A,
\]
where $H$ is a hyperbolic plane and $A$ is anisotropic. Since
$\dim V=\dim U$, we have
\[
2r+\dim A=2s+\dim A,
\]
so $r=s$. Therefore $(V,\varphi)\cong (U,\psi)$.
\end{solution}


\subsection*{Exercise A.11.3}

Let $K$ be a field of characteristic $\neq 2$.

\begin{enumerate}
    \item Prove that for any two nonsingular quadratic spaces $(U,\psi)$, $(V,\varphi)$ over $K$, one has
    \[
        \operatorname{disc}\bigl((U,\psi)\perp (V,\varphi)\bigr)
        =
        (-1)^{\dim U \dim V}
        \operatorname{disc}(U,\psi)\operatorname{disc}(V,\varphi)
        \in K^{\times}/K^{\times 2}.
    \]

    \item Let $(V,\varphi)$ be a nonsingular quadratic space of dimension $2$ over $K$. Prove that $(V,\varphi)$ is hyperbolic if and only if $\operatorname{disc}(V,\varphi)=1\in K^{\times}/K^{\times 2}$.
\end{enumerate}

\begin{solution}
    We use the convention that, if $(V,\varphi)$ has dimension $n$ and Gram matrix $A$, then \[ \operatorname{disc}(V,\varphi) = (-1)^{n(n-1)/2}\det(A) \in K^\times/K^{\times 2}. \] For the first statement, let $\dim U=m$ and $\dim V=n$. Choose Gram matrices $A$ and $B$ for $(U,\psi)$ and $(V,\varphi)$ respectively. Then the Gram matrix of $(U,\psi)\perp (V,\varphi)$ is \[ \begin{pmatrix} A & 0\\ 0 & B \end{pmatrix}. \] Hence \[ \operatorname{disc}\bigl((U,\psi)\perp (V,\varphi)\bigr) = (-1)^{(m+n)(m+n-1)/2}\det(A)\det(B). \] On the other hand, \[ \operatorname{disc}(U,\psi)\operatorname{disc}(V,\varphi) = (-1)^{m(m-1)/2+n(n-1)/2}\det(A)\det(B). \] Since \[ \frac{(m+n)(m+n-1)}{2} - \frac{m(m-1)}{2} - \frac{n(n-1)}{2} = mn, \] we get \[ \operatorname{disc}\bigl((U,\psi)\perp (V,\varphi)\bigr) = (-1)^{mn} \operatorname{disc}(U,\psi)\operatorname{disc}(V,\varphi). \] For the second statement, since $\operatorname{char}K\neq 2$, we may write \[ (V,\varphi)\cong \langle a,b\rangle \] with $a,b\in K^\times$. Then \[ \operatorname{disc}(V,\varphi)=-ab\in K^\times/K^{\times 2}. \] If $(V,\varphi)$ is hyperbolic, then $(V,\varphi)\cong \langle 1,-1\rangle$, so \[ \operatorname{disc}(V,\varphi)=-(1)(-1)=1. \] Conversely, suppose that $\operatorname{disc}(V,\varphi)=1$. Then $-ab\in K^{\times 2}$, so there exists $c\in K^\times$ such that $-ab=c^2$. Thus \[ b=-\frac{c^2}{a}=-a\left(\frac{c}{a}\right)^2. \] Therefore \[ \langle a,b\rangle = \left\langle a,-\frac{c^2}{a}\right\rangle \cong \langle a,-a\rangle. \] The vector $(1,1)$ is isotropic for $\langle a,-a\rangle$. Since $\langle a,-a\rangle$ is nonsingular of dimension $2$, it is hyperbolic. Hence $(V,\varphi)$ is hyperbolic.
\end{solution}

\subsection*{Exercise A.11.4}

Let $F$ be a henselian discrete valuation field with residue field $k$. Assume $\operatorname{char}(k)\neq 2$. Let $u_1,\cdots,u_r,w_1,\cdots,w_r\in \mathcal{O}_F^{\times}$, and put
\[
    \varphi=\langle u_1,\cdots,u_r\rangle,
    \qquad
    \psi=\langle w_1,\cdots,w_r\rangle.
\]

\begin{enumerate}
    \item Prove that $\varphi\perp -\psi$ is hyperbolic over $F$ if and only if $\overline{\varphi}\perp -\overline{\psi}$ is hyperbolic over $k$.

    \item Prove that $\varphi$ and $\psi$ are isomorphic over $F$ if and only if $\overline{\varphi}$ and $\overline{\psi}$ are isomorphic over $k$.
\end{enumerate}
\begin{solution}
    We use the following consequence of Springer's theorem: if $q=\langle a_1,\dots,a_n\rangle$ with all $a_i\in \mathcal O_F^\times$, then $q$ is hyperbolic over $F$ if and only if its reduction $\overline q=\langle \overline a_1,\dots,\overline a_n\rangle$ is hyperbolic over $k$. Equivalently, the reduction map on unit forms detects the zero class in the Witt group. Now \[ \varphi\perp -\psi = \langle u_1,\dots,u_r,-w_1,\dots,-w_r\rangle \] is a unit form over $F$, and its reduction is \[ \overline\varphi\perp -\overline\psi = \langle \overline u_1,\dots,\overline u_r, -\overline w_1,\dots,-\overline w_r\rangle . \] Therefore, by Springer's theorem, \[ \varphi\perp -\psi \text{ is hyperbolic over } F \quad\Longleftrightarrow\quad \overline\varphi\perp -\overline\psi \text{ is hyperbolic over } k. \] This proves the first statement. For the second statement, since $\varphi$ and $\psi$ have the same dimension, Exercise A.11.2(3) gives \[ \varphi\cong \psi \quad\Longleftrightarrow\quad \varphi\perp -\psi \text{ is hyperbolic over } F. \] Similarly, \[ \overline\varphi\cong \overline\psi \quad\Longleftrightarrow\quad \overline\varphi\perp -\overline\psi \text{ is hyperbolic over } k. \] By the first statement, the two hyperbolicity conditions are equivalent. Hence \[ \varphi\cong \psi \quad\Longleftrightarrow\quad \overline\varphi\cong \overline\psi . \]
\end{solution}



\subsection*{Exercise A.11.5}

Let $\mathbb{F}$ be a finite field of characteristic $\neq 2$. Prove that up to isomorphism there is one and only one anisotropic binary quadratic space over $\mathbb{F}$.

\begin{solution}

Let $\delta\in \mathbb F^\times$ be a nonsquare. We claim that
\[
\langle 1,-\delta\rangle
\]
is the unique anisotropic binary quadratic space over $\mathbb F$ up to
isomorphism.

First, $\langle 1,-\delta\rangle$ is anisotropic. Indeed, if
$x^2-\delta y^2=0$ and $y\neq 0$, then $\delta=(x/y)^2$, contradicting that
$\delta$ is a nonsquare. Hence $y=0$, and then $x=0$.

Now let $(V,q)$ be any anisotropic binary quadratic space over $\mathbb F$.
Since $\operatorname{char}\mathbb F\neq 2$, we may diagonalize
\[
(V,q)\cong \langle a,b\rangle
\]
with $a,b\in \mathbb F^\times$. By Exercise A.11.3(2), a nonsingular binary
form is hyperbolic if and only if its discriminant is $1$. Since $(V,q)$ is
anisotropic, it is not hyperbolic, so
\[
\operatorname{disc}(V,q)=-ab
\]
is the nontrivial square class in $\mathbb F^\times/\mathbb F^{\times 2}$.
Thus $-ab=\delta$ in $\mathbb F^\times/\mathbb F^{\times 2}$, and hence
\[
\langle a,b\rangle \cong \left\langle a,-\frac{\delta}{a}\right\rangle .
\]

It remains to compare this form with $\langle 1,-\delta\rangle$. Since
$\mathbb F(\sqrt{\delta})/\mathbb F$ is a quadratic extension of finite fields,
its norm map is surjective. Hence the norm form
\[
x^2-\delta y^2
\]
represents $a$. Therefore
\[
\langle 1,-\delta\rangle \cong \langle a,c\rangle
\]
for some $c\in\mathbb F^\times$. Comparing discriminants gives
\[
-ac=\delta \quad \text{in } \mathbb F^\times/\mathbb F^{\times 2},
\]
so $c$ differs from $-\delta/a$ by a square factor. Hence
\[
\langle a,c\rangle \cong \left\langle a,-\frac{\delta}{a}\right\rangle .
\]
Therefore
\[
(V,q)\cong \langle 1,-\delta\rangle .
\]
Thus, up to isomorphism, there is exactly one anisotropic binary quadratic
space over $\mathbb F$.
\end{solution}

\subsection*{Exercise A.11.6}

Let $F$ be a henselian discrete valuation field of characteristic $\neq 2$. Prove that $F^{\times 2}$ is an open and closed subgroup of $F^{\times}$. Here of course, the topology on $F^{\times}$ is the subspace topology induced from $F$.

\begin{solution}
    Let $v$ be the discrete valuation on $F$, let $\mathcal O=\mathcal O_F$, and let $\mathfrak m$ be its maximal ideal. Put $e=v(2)$, and choose an integer $N>2e$. We first prove that \[ 1+\mathfrak m^N\subseteq F^{\times 2}. \] Let $u\in 1+\mathfrak m^N$. If $u=1$, there is nothing to prove. Otherwise write $u=1+t$ with $v(t)\geq N$. We want to solve $X^2=u$. Put $X=1+Y$. Then this is equivalent to \[ Y^2+2Y-t=0. \] Consider \[ g(Y)=Y^2+2Y-t. \] The relevant points of its Newton polygon are \[ (0,v(t)),\qquad (1,e),\qquad (2,0). \] Since $v(t)\geq N>2e$, the point $(1,e)$ lies strictly below the line joining $(0,v(t))$ and $(2,0)$. Hence the Newton polygon has two sides, \[ (0,v(t))\to (1,e),\qquad (1,e)\to (2,0). \] The first side has horizontal length $1$ and slope $e-v(t)$. By the Newton polygon theorem over a henselian discrete valuation field, $g$ has a root $Y\in F$ with \[ v(Y)=v(t)-e>0. \] Thus $X=1+Y\in F^\times$ satisfies $X^2=u$. Therefore \[ 1+\mathfrak m^N\subseteq F^{\times 2}. \] It follows that $F^{\times 2}$ is open in $F^\times$. Indeed, for every $a^2\in F^{\times 2}$, the set \[ a^2(1+\mathfrak m^N) \] is an open neighborhood of $a^2$ contained in $F^{\times 2}$. Finally, $F^{\times 2}$ is a subgroup of the topological group $F^\times$. Every coset of an open subgroup is open, so the complement of $F^{\times 2}$ is a union of open cosets. Hence $F^{\times 2}$ is closed. Therefore $F^{\times 2}$ is an open and closed subgroup of $F^\times$.
\end{solution}


\section*{Homework 12}

In this section, let $F$ be a henselian discrete valuation field of
characteristic $\neq 2$, let $\pi\in \mathcal O_F$ be a uniformizer and
$k=\mathcal O_F/\pi\mathcal O_F$. The normalized discrete valuation on $F$ is
denoted by $v$ or $v_F$.

\subsection*{Exercise A.12.1} Let $\mathfrak a$ be a fractional ideal of $F$. 
\begin{enumerate} 
    \item For any $\alpha,\beta\in F^\times$, let us write $\alpha\simeq \beta$ if $\alpha\beta^{-1}\in \mathcal O_F^{\times 2}$. 
        \begin{enumerate} 
            \item Prove that the above relation $\simeq$ is an equivalence relation on the set $F^\times$. 
            \item Show that for any $\alpha,\beta\in F^\times$, we have \[ \alpha\simeq \beta \Longleftrightarrow \alpha\beta^{-1}\in \mathcal O_F^\times \text{ and } \mathfrak d(\alpha\beta^{-1})=0. \] 
        \end{enumerate} 
        \item For any $\alpha,\beta\in F^\times$, let us write $\alpha\simeq \beta \pmod{\mathfrak a}$ if $\alpha\beta^{-1}\in \mathcal O_F^\times$ and \[ \alpha=\beta\varepsilon^2 \pmod{\mathfrak a} \] for some $\varepsilon\in \mathcal O_F^\times$. 
            \begin{enumerate} 
                \item Prove that the above congruence relation modulo $\mathfrak a$ is an equivalence relation on $F^\times$. 
                \item Show that if $\alpha,\beta\in F^\times$ satisfies $\alpha\simeq \beta\pmod{\mathfrak a}$, then \[ \mathfrak d(\alpha\beta^{-1})\subseteq \beta^{-1}\mathfrak a. \] 
            \end{enumerate} 
        \item Assume $k$ is perfect of characteristic $2$. Let $\alpha,\beta\in F^\times$. 
            \begin{enumerate} 
                \item Prove that the relation $\alpha\simeq \beta\pmod{\mathfrak a}$ holds if and only if $\alpha\beta^{-1}\in \mathcal O_F^\times$ and \[ \mathfrak d(\alpha\beta^{-1})\subseteq \beta^{-1}\mathfrak a. \] 
                \item Show that if $\alpha\beta^{-1}\in \mathcal O_F^\times$, then \[ \alpha\simeq \beta \pmod{\alpha\pi\mathcal O_F}. \] 
            \end{enumerate} 
\end{enumerate}
\begin{solution}

Write $\mathcal O=\mathcal O_F$. We use the following description of the
quadratic defect:
\[
    \mathfrak d(c)=\bigcap_{x\in F}(c-x^2)\mathcal O .
\]
Thus $\mathfrak d(c)\subseteq (c-x^2)\mathcal O$ for every $x\in F$.

\paragraph*{(1.a)}
The relation $\simeq$ is the equivalence relation induced by the subgroup
$\mathcal O^{\times 2}\subseteq F^\times$. Indeed,
$\alpha\alpha^{-1}=1\in\mathcal O^{\times 2}$, if
$\alpha\beta^{-1}\in\mathcal O^{\times 2}$ then
$\beta\alpha^{-1}=(\alpha\beta^{-1})^{-1}\in\mathcal O^{\times 2}$, and if
$\alpha\beta^{-1},\beta\gamma^{-1}\in\mathcal O^{\times 2}$, then
$\alpha\gamma^{-1}=(\alpha\beta^{-1})(\beta\gamma^{-1})\in
\mathcal O^{\times 2}$.

\paragraph*{(1.b)}
Put $c=\alpha\beta^{-1}$. If $\alpha\simeq\beta$, then
$c\in\mathcal O^{\times 2}$, hence $c\in\mathcal O^\times$ and
$\mathfrak d(c)=0$.

Conversely, suppose $c\in\mathcal O^\times$ and $\mathfrak d(c)=0$. Then $c$
can be approximated arbitrarily closely by squares. Choose $x\in F$ such that
$v(c-x^2)>2v(2)$. Since $c\in\mathcal O^\times$, this forces $x\in\mathcal
O^\times$. Applying Hensel's lemma to $T^2-c$ at $x$, we get $c\in F^{\times
2}$. Since $v(c)=0$, its square root is a unit. Hence
$c\in\mathcal O^{\times 2}$, so $\alpha\simeq\beta$.

\paragraph*{(2.a)}
Reflexivity is clear by taking $\varepsilon=1$. For symmetry, suppose
$\alpha=\beta\varepsilon^2 \pmod{\mathfrak a}$. Then
\[
    \beta=\alpha\varepsilon^{-2}\pmod{\mathfrak a},
\]
because $\varepsilon^{-2}\in\mathcal O^\times$ and $\mathfrak a$ is an
$\mathcal O$-submodule of $F$. Also $\beta\alpha^{-1}\in\mathcal O^\times$.
Hence $\beta\simeq\alpha\pmod{\mathfrak a}$.

For transitivity, suppose
\[
    \alpha=\beta\varepsilon^2\pmod{\mathfrak a},
    \qquad
    \beta=\gamma\eta^2\pmod{\mathfrak a}.
\]
Then
\[
    \alpha=\gamma(\eta\varepsilon)^2\pmod{\mathfrak a},
\]
and $\alpha\gamma^{-1}\in\mathcal O^\times$. Hence
$\alpha\simeq\gamma\pmod{\mathfrak a}$.

\paragraph*{(2.b)}
Put $c=\alpha\beta^{-1}$. Since
$\alpha\simeq\beta\pmod{\mathfrak a}$, there exists
$\varepsilon\in\mathcal O^\times$ such that
\[
    \alpha-\beta\varepsilon^2\in\mathfrak a.
\]
Equivalently,
\[
    c-\varepsilon^2\in\beta^{-1}\mathfrak a.
\]
By the definition of quadratic defect,
\[
    \mathfrak d(c)\subseteq (c-\varepsilon^2)\mathcal O
    \subseteq \beta^{-1}\mathfrak a.
\]
Thus
\[
    \mathfrak d(\alpha\beta^{-1})\subseteq\beta^{-1}\mathfrak a.
\]

\paragraph*{(3.a)}
Put $c=\alpha\beta^{-1}$. The implication
\[
    \alpha\simeq\beta\pmod{\mathfrak a}
    \Longrightarrow
    c\in\mathcal O^\times
    \text{ and }
    \mathfrak d(c)\subseteq\beta^{-1}\mathfrak a
\]
is exactly (2.b).

Conversely, suppose $c\in\mathcal O^\times$ and
$\mathfrak d(c)\subseteq\beta^{-1}\mathfrak a$. Since $k$ is perfect of
characteristic $2$, every element of $k^\times$ is a square. Hence $c$ is
congruent modulo $\pi\mathcal O$ to a unit square. Therefore, by discreteness
of the valuation, there exists $\varepsilon\in\mathcal O^\times$ such that
\[
    \mathfrak d(c)=(c-\varepsilon^2)\mathcal O .
\]
Thus
\[
    c-\varepsilon^2\in \mathfrak d(c)\subseteq \beta^{-1}\mathfrak a.
\]
Multiplying by $\beta$, we get
\[
    \alpha-\beta\varepsilon^2\in\mathfrak a.
\]
Hence
\[
    \alpha\simeq\beta\pmod{\mathfrak a}.
\]

\paragraph*{(3.b)}
Put $c=\alpha\beta^{-1}\in\mathcal O^\times$. Since $k$ is perfect of
characteristic $2$, choose $\varepsilon\in\mathcal O^\times$ such that
\[
    c\equiv \varepsilon^2\pmod{\pi\mathcal O}.
\]
Then
\[
    c-\varepsilon^2\in\pi\mathcal O=c\pi\mathcal O,
\]
because $c$ is a unit. Multiplying by $\beta$, we obtain
\[
    \alpha-\beta\varepsilon^2\in \beta c\pi\mathcal O
    =
    \alpha\pi\mathcal O.
\]
Therefore
\[
    \alpha\simeq\beta\pmod{\alpha\pi\mathcal O}.
\]
\end{solution}


\subsection*{Exercise A.12.2}

For a nonempty subset $S\subseteq F$, define
\[
    \operatorname{ord}(S):=
    \inf\{v(x)\mid x\in S\}\in \mathbb Z\cup\{\pm\infty\}.
\]

For any $a\in F^\times$ we define its \emph{relative quadratic defect} as
\[
    \operatorname{dft}(a):=
    \operatorname{ord}\bigl(a^{-1}\mathfrak d(a)\bigr)
    \in \mathbb Z\cup\{+\infty\}.
\]

Prove that $\operatorname{dft}(a)\in \mathbb N\cup\{+\infty\}$ for all
$a\in F^\times$, and that there is a well-defined map
\[
    \operatorname{dft}:F^\times/F^{\times 2}
    \longrightarrow
    \mathbb N\cup\{+\infty\},
    \qquad
    aF^{\times 2}\longmapsto \operatorname{dft}(a).
\]

Let $a\in F^\times$ and put $\varepsilon=a\pi^{-v(a)}$. Prove
\[
    \operatorname{dft}(a)=
    \begin{cases}
        0, & \text{if } v(a)=R \text{ is odd},\\
        \operatorname{dft}(\varepsilon)
        =\operatorname{ord}\bigl(\mathfrak d(\varepsilon)\bigr),
        & \text{if } v(a)=R \text{ is even}.
    \end{cases}
\]

\begin{solution}
We use the definition
\[
    \mathfrak d(a)=\bigcap_{x\in F}(a-x^2)\mathcal O_F .
\]
Since the term with $x=0$ occurs in the intersection, we have
\[
    \mathfrak d(a)\subseteq a\mathcal O_F.
\]
Hence
\[
    a^{-1}\mathfrak d(a)\subseteq \mathcal O_F,
\]
so
\[
    \operatorname{dft}(a)
    =
    \operatorname{ord}\bigl(a^{-1}\mathfrak d(a)\bigr)
    \in \mathbb N\cup\{+\infty\}.
\]

We next show that $\operatorname{dft}(a)$ only depends on the square class of
$a$. Let $b\in F^\times$. Then
\[
    \mathfrak d(ab^2)
    =
    \bigcap_{x\in F}(ab^2-x^2)\mathcal O_F.
\]
Writing $x=by$, we get
\[
    ab^2-x^2=b^2(a-y^2).
\]
Since $y$ also runs through all of $F$, it follows that
\[
    \mathfrak d(ab^2)=b^2\mathfrak d(a).
\]
Therefore
\[
    (ab^2)^{-1}\mathfrak d(ab^2)
    =
    a^{-1}\mathfrak d(a),
\]
and hence
\[
    \operatorname{dft}(ab^2)=\operatorname{dft}(a).
\]
Thus
\[
    \operatorname{dft}:F^\times/F^{\times 2}
    \longrightarrow
    \mathbb N\cup\{+\infty\}
\]
is well-defined.

Now let $a\in F^\times$ and put
\[
    R=v(a),\qquad \varepsilon=a\pi^{-R}.
\]
Then $\varepsilon\in\mathcal O_F^\times$.

If $R$ is odd, then for every $x\in F$, the two valuations $v(a)=R$ and
$v(x^2)=2v(x)$ are distinct. Hence
\[
    v(a-x^2)=\min\{R,2v(x)\}\leq R.
\]
Taking $x=0$ gives equality, so
\[
    \operatorname{ord}(\mathfrak d(a))=R.
\]
Therefore
\[
    \operatorname{dft}(a)
    =
    \operatorname{ord}\bigl(a^{-1}\mathfrak d(a)\bigr)
    =
    R-R
    =
    0.
\]

If $R$ is even, write $R=2s$. Then
\[
    a=\pi^{2s}\varepsilon.
\]
For $x\in F$, write $x=\pi^s y$. Then
\[
    a-x^2
    =
    \pi^{2s}(\varepsilon-y^2).
\]
Since $y$ runs through all of $F$, we get
\[
    \mathfrak d(a)=\pi^{2s}\mathfrak d(\varepsilon).
\]
Hence
\[
    a^{-1}\mathfrak d(a)
    =
    \varepsilon^{-1}\mathfrak d(\varepsilon).
\]
Because $\varepsilon$ is a unit,
\[
    \operatorname{dft}(a)
    =
    \operatorname{dft}(\varepsilon)
    =
    \operatorname{ord}\bigl(\mathfrak d(\varepsilon)\bigr).
\]
Thus
\[
    \operatorname{dft}(a)=
    \begin{cases}
        0, & \text{if } v(a) \text{ is odd},\\
        \operatorname{dft}(\varepsilon)
        =
        \operatorname{ord}\bigl(\mathfrak d(\varepsilon)\bigr),
        & \text{if } v(a) \text{ is even}.
    \end{cases}
\]
\end{solution}

\subsection*{Exercise A.12.3}

With notation as in Exercise A.12.2, now assume further that the residue field
$k$ is a perfect field of characteristic $2$ which has a unique quadratic
extension. Let $\Delta\in \mathcal O_F^\times$ be chosen such that
\[
    \mathfrak d(\Delta)=4\mathcal O_F
\]
(cf. Lemma 4.2.12). Let $e=v_F(2)$.

Prove
\[
    \operatorname{dft}(F^\times)
    =
    \{0,2e,+\infty\}
    \cup
    \{2j-1\mid j\in [1,e]\}.
\]

Show that for any $a\in F^\times$, we have the following equivalences:
\[
    \operatorname{dft}(a)=0
    \Longleftrightarrow
    v(a)\text{ is odd},
\]
\[
    \operatorname{dft}(a)\geq 1
    \Longleftrightarrow
    v(a)\text{ is even},
\]
\[
    \operatorname{dft}(a)=2e
    \Longleftrightarrow
    aF^{\times 2}=\Delta F^{\times 2},
\]
\[
    \operatorname{dft}(a)=+\infty
    \Longleftrightarrow
    a\in F^{\times 2}.
\]

For any $a,b\in F^\times$, prove that
\[
    \begin{cases}
        \operatorname{dft}(ab)=0
        =
        \min\{\operatorname{dft}(a),\operatorname{dft}(b)\},
        & \text{if } v(a)\not\equiv v(b)\pmod 2,\\
        \operatorname{dft}(ab)=0
        =
        \operatorname{dft}(a)=\operatorname{dft}(b),
        & \text{if } v(a)\equiv v(b)\equiv 1\pmod 2,
    \end{cases}
\]
and that we have the \emph{domination principle}
\[
    \operatorname{dft}(ab)
    \geq
    \min\{\operatorname{dft}(a),\operatorname{dft}(b)\},
\]
and if $\operatorname{dft}(a)\neq \operatorname{dft}(b)$, then
\[
    \operatorname{dft}(ab)
    =
    \min\{\operatorname{dft}(a),\operatorname{dft}(b)\}.
\]
\begin{solution}
    Let $\mathcal O=\mathcal O_F$ and let $e=v(2)$. We use the following standard consequence of Lemma 4.2.12. For $u\in \mathcal O^\times$, \[ \operatorname{dft}(u)\in \{1,3,\dots,2e-1,2e,+\infty\}. \] Moreover, \[ \operatorname{dft}(u)=+\infty \Longleftrightarrow u\in F^{\times 2}, \] and \[ \operatorname{dft}(u)=2e \Longleftrightarrow uF^{\times 2}=\Delta F^{\times 2}. \] The values $1,3,\dots,2e-1$ occur, for instance, on suitable unit classes represented by elements of the form $1+c\pi^{2j-1}$. By Exercise A.12.2, if $v(a)$ is odd, then \[ \operatorname{dft}(a)=0. \] If $v(a)$ is even, write $a=\pi^{2s}u$ with $u\in\mathcal O^\times$. Since $\pi^{2s}$ is a square, $a$ and $u$ have the same square class, and \[ \operatorname{dft}(a)=\operatorname{dft}(u). \] Therefore \[ \operatorname{dft}(F^\times) = \{0,2e,+\infty\} \cup \{2j-1\mid 1\leq j\leq e\}. \] Next we prove the stated equivalences. The first two follow immediately from Exercise A.12.2 and the unit classification above: \[ \operatorname{dft}(a)=0 \Longleftrightarrow v(a)\text{ is odd}, \] and \[ \operatorname{dft}(a)\geq 1 \Longleftrightarrow v(a)\text{ is even}. \] If $v(a)$ is even, then $aF^{\times 2}=uF^{\times 2}$ for some $u\in\mathcal O^\times$. Hence \[ \operatorname{dft}(a)=2e \Longleftrightarrow uF^{\times 2}=\Delta F^{\times 2} \Longleftrightarrow aF^{\times 2}=\Delta F^{\times 2}. \] Similarly, \[ \operatorname{dft}(a)=+\infty \Longleftrightarrow a\in F^{\times 2}. \] It remains to prove the domination principle. Since $\operatorname{dft}$ only depends on square classes, we may multiply $a$ and $b$ by squares whenever convenient. If $v(a)\not\equiv v(b)\pmod 2$, then $v(ab)$ is odd. Hence, by Exercise A.12.2, \[ \operatorname{dft}(ab)=0 = \min\{\operatorname{dft}(a),\operatorname{dft}(b)\}. \] If $v(a)\equiv v(b)\equiv 1\pmod 2$, then \[ \operatorname{dft}(a)=\operatorname{dft}(b)=0. \] In this case one only gets $\operatorname{dft}(ab)\geq 0$ in general. Now assume $v(a)$ and $v(b)$ are both even. After multiplying by squares, we may assume $a,b\in\mathcal O^\times$. Put \[ r=\operatorname{dft}(a), \qquad s=\operatorname{dft}(b), \qquad m=\min\{r,s\}. \] If one of $r,s$ is $+\infty$, the assertion is immediate from the fact that squares do not change square classes. Thus assume $r,s<+\infty$. By Exercise A.12.1(3.a), there exist $\varepsilon,\eta\in\mathcal O^\times$ such that \[ a\equiv \varepsilon^2 \pmod{a\pi^r\mathcal O}, \qquad b\equiv \eta^2 \pmod{b\pi^s\mathcal O}. \] Since $a,b,\varepsilon,\eta$ are units, multiplying the two congruences gives \[ ab\equiv (\varepsilon\eta)^2 \pmod{ab\pi^m\mathcal O}. \] Again by Exercise A.12.1(3.a), \[ \operatorname{dft}(ab)\geq m. \] Thus \[ \operatorname{dft}(ab) \geq \min\{\operatorname{dft}(a),\operatorname{dft}(b)\}. \] Finally assume $\operatorname{dft}(a)\neq\operatorname{dft}(b)$. Without loss of generality, let $r<s$. The domination principle gives \[ \operatorname{dft}(ab)\geq r. \] If $\operatorname{dft}(ab)>r$, then applying the domination principle to $a=(ab)b^{-1}$ gives \[ r=\operatorname{dft}(a) = \operatorname{dft}((ab)b^{-1}) \geq \min\{\operatorname{dft}(ab),\operatorname{dft}(b^{-1})\}. \] Since $\operatorname{dft}(b^{-1})=\operatorname{dft}(b)=s$ and both $\operatorname{dft}(ab)$ and $s$ are strictly larger than $r$, this is a contradiction. Hence \[ \operatorname{dft}(ab)=r = \min\{\operatorname{dft}(a),\operatorname{dft}(b)\}. \]
\end{solution}

\end{document}